%% Volume I, Chapter 23: First-Order Equations
\chapter{First-Order Equations}
\label{ch:first-order-equations}

\abstract{Many laws of nature say how a quantity changes rather than what it is. A radioactive sample loses a fixed fraction of itself per unit of time, a cup of coffee cools in proportion to its excess temperature, a population grows until resources run short. Each such law is a differential equation, and to solve it is to recover the whole history from the local rule and one observation. Few equations can be solved by formulas; this chapter collects the methods for equations of the first order. For the others three questions remain, and they run through this part of the book: whether a solution exists and is unique, what it looks like, and how a computer can follow it.}

Differential equations have appeared several times in this book already. \Cref{thm:exp-ode} showed that the solutions of $f'=\lambda f$ are exactly the functions $ce^{\lambda x}$, and \cref{lem:harmonic-unique} that those of $f''=-f$ are the functions $a\cos x+b\sin x$; both results characterised elementary functions by the equations they satisfy. \Cref{thm:linear-system} extended the first of them to systems $\mathbf x'=A\mathbf x$, whose solutions are $e^{tA}\mathbf x_0$. Liouville's criterion for elementary antiderivatives, on which the proof of \cref{thm:liouville} rested, was a statement about the rational solutions of the equation $R'+g'R=h$, and the catenary of \cref{sec:hyperbolic} was described by the equation $y''=\sqrt{1+y'^2}/a$, which was left unsolved. This part of the book studies such equations systematically. The present chapter treats equations of the first order, \cref{ch:linear-higher-order} linear equations of higher order, Chapter~25 linear systems, Chapter~26 solutions by power series and Chapter~27 the discrete counterparts of differential equations, the difference equations and recurrences.

A first-order equation $y'=f(t,y)$ prescribes the rate of change of an unknown function $y$ in terms of the time $t$ and of the current value $y(t)$. Three kinds of answers can be given. For special forms of $f$ the solutions can be written down: the equations whose variables separate, the linear equations, the exact equations and those that a substitution brings into one of these forms; \cref{sec:separable,sec:linear-first-order,sec:exact,sec:substitutions} develop these methods and apply them to growth and decay, to cooling, to populations, to falling bodies and to the hanging chain. For most $f$ no formula exists, and then the direction field of \cref{sec:direction-fields} and the qualitative theory of \cref{sec:autonomous} describe the solutions without computing them. Underneath both lies the theory of \cref{sec:existence-uniqueness}, which says when a solution exists, when it is unique and how it depends on its initial value; the examples of \cref{sec:ode-basics} show that none of this can be taken for granted. Finally, \cref{sec:euler} computes solutions approximately with the simplest numerical method, that of Euler, proves that it converges, and observes its limitations in experiments.

We fix the notation for the whole of Part~IV. The independent variable is $t$, which is often, though not always, a time, and the unknown function is $y$, or $\mathbf x$ for a vector-valued unknown, as in \cref{thm:linear-system}. Derivatives are written with primes, $y'=dy/dt$ and $y''=d^2y/dt^2$; we do not use the dots $\dot y$, $\ddot y$ with which texts on mechanics often denote derivatives with respect to time. The same letter denotes a function and its value, so that $y'=f(t,y)$ abbreviates $y'(t)=f(t,y(t))$. In problems about curves in the plane, such as the exact equations of \cref{sec:exact} and the catenary, where a solution is a curve rather than a history, we use the coordinates $x$ and $y$ and write $y'=dy/dx$, as in the earlier chapters.

The chapter is organised as follows. \Cref{sec:ode-basics} introduces the vocabulary, the notions of solution and of initial value problem, two warning examples and the equivalent integral equation. \Cref{sec:direction-fields} draws direction fields. \Cref{sec:separable,sec:linear-first-order,sec:exact,sec:substitutions} treat separable, linear and exact equations and substitutions, \cref{sec:autonomous} autonomous equations and their stability, and \cref{sec:existence-uniqueness} existence, uniqueness and continuous dependence. \Cref{sec:euler} is a first look at numerical solution.

\section{Differential Equations and Their Solutions}
\label{sec:ode-basics}

\subsection{Terminology}

A \term{differential equation} is an equation that relates an unknown function to its derivatives. If the unknown function depends on one variable, the equation is an \term{ordinary differential equation}; if it depends on several variables and the equation involves its partial derivatives, it is a \term{partial differential equation}, such as the heat equation $u_t=u_{xx}$ for a temperature $u(x,t)$. This volume treats ordinary differential equations only. The \term{order} of an equation is the order of the highest derivative that occurs in it. The equations
\begin{equation}
y'=ky,\qquad y'=t-y,\qquad y'=y(1-y),\qquad y'=t^2+y^2,\qquad y''+y=0
\label{eq:first-examples}
\end{equation}
are ordinary differential equations, the first four of the first order and the last of the second. An equation of order $k$ is \term{linear} if the unknown function and its derivatives occur only linearly, that is, if it has the form
\[
a_k(t)y^{(k)}+a_{k-1}(t)y^{(k-1)}+\cdots+a_1(t)y'+a_0(t)y=q(t)
\]
with given functions $a_0,\ldots,a_k$ and $q$; otherwise it is \term{nonlinear}. In \eqref{eq:first-examples} the first, the second and the last equation are linear, and the third and the fourth are nonlinear, because of the terms $y^2$. An equation is \term{autonomous} if the independent variable does not occur in it explicitly: $y'=ky$, $y'=y(1-y)$ and $y''+y=0$ are autonomous, while in $y'=t-y$ and $y'=t^2+y^2$ the rate of change depends on the time as well as on the state.

This chapter studies equations of the first order solved for the derivative,
\begin{equation}
y'=f(t,y),
\label{eq:first-order}
\end{equation}
with a given function $f$ of two variables. Several unknown functions satisfying as many equations form a \term{system} of differential equations, such as the system $\mathbf x'=A\mathbf x$ of \cref{thm:linear-system}. An equation of higher order is equivalent to a system of the first order: the second-order equation $y''=g(t,y,y')$ holds exactly when the pair $(y,v)$ with $v\coloneqq y'$ satisfies the system $y'=v$, $v'=g(t,y,v)$. Chapter~25 uses this reduction; in this chapter it will occasionally turn a second-order equation into a first-order equation for $v$ alone (\cref{sec:substitutions}).

\subsection{Solutions and Initial Value Problems}

\begin{definition}[Solution, initial value problem]
\label{def:ode-solution}
Let $D\subseteq\mathbb R^2$ and $f\colon D\to\mathbb R$. A \term{solution} of the equation $y'=f(t,y)$ on an interval $I$ containing more than one point is a differentiable function $y\colon I\to\mathbb R$ with $(t,y(t))\in D$ and
\[
y'(t)=f\bigl(t,y(t)\bigr)\qquad\text{for all }t\in I ,
\]
the derivative being one-sided at an end of $I$ that belongs to $I$. For $(t_0,y_0)\in D$, a solution of the \term{initial value problem}
\begin{equation}
y'=f(t,y),\qquad y(t_0)=y_0
\label{eq:ivp}
\end{equation}
is a solution on an interval $I$ containing $t_0$ with $y(t_0)=y_0$. The number $y_0$ is the \term{initial value}, and $(t_0,y_0)$ the \term{initial point}.
\end{definition}

A solution is a function together with its interval of definition. The function $y(t)=1/(1-t)$ satisfies $y'=1/(1-t)^2=y^2$ wherever it is defined, but its domain $\mathbb R\setminus\{1\}$ is not an interval; it gives one solution on $(-\infty,1)$ and another on $(1,\infty)$, which have nothing to do with each other, since no solution on an interval can pass from one to the other. The insistence on intervals is not pedantry: on a set with a gap the derivative does not determine the function, by the remark after \cref{cor:zero-derivative}.

The simplest equations have a right side that does not depend on $y$. For a continuous $f\colon J\to\mathbb R$ on an interval $J$, the solutions of $y'=f(t)$ are the antiderivatives of $f$, and by \cref{cor:antiderivative} the initial value problem $y'=f(t)$, $y(t_0)=y_0$, has exactly one solution on $J$,
\[
y(t)=y_0+\int_{t_0}^tf(s)\,ds .
\]
This is why solving a differential equation is traditionally called \emph{integrating} it, and why a formula for the solutions that uses finitely many integrations of known functions is called a solution \term{by quadratures}. The example already shows two features of the general case. The solutions form a family with one free constant, here the additive constant of the antiderivative; such a family is called a \term{general solution}, and each member a \term{particular solution}. And one additional condition, the initial value, singles out one member of the family.

The equation $y'=ky$ with a constant $k$ is the next simplest. By \cref{thm:exp-ode}, applied after a shift of the origin of time when the interval does not contain $0$, its solutions on any interval are the functions $y(t)=ce^{kt}$, and the initial value problem with $y(t_0)=y_0$ has the unique solution $y_0e^{k(t-t_0)}$, defined for all $t$. Here the free constant enters as a factor.

\subsection{Equations from Models}

Differential equations arise when a law describes the rate of change of a quantity. The following laws will serve as examples throughout the chapter; in each, the equation is an idealisation, and the constants must be determined by observation.

The amount $N(t)$ of a radioactive substance decreases at a rate proportional to the amount present, $N'=-\lambda N$ with a constant $\lambda>0$, as \cref{sec:exp-log-numerics} recalled. A body at temperature $T(t)$ in surroundings at the constant temperature $A$ loses heat at a rate roughly proportional to the difference of the temperatures, which gives Newton's law of cooling,
\[
T'=-k(T-A),\qquad k>0 .
\]
A body of mass $m$ falling through air feels its weight $mg$, with the acceleration of gravity $g$, and a resistance that grows with its speed $v$. For slow motion in a viscous medium the resistance is roughly proportional to $v$, and for fast motion in air roughly proportional to $v^2$, so that Newton's second law, mass times acceleration equals force, gives
\[
mv'=mg-kv\qquad\text{or}\qquad mv'=mg-kv^2,\qquad k>0,
\]
with the downward direction taken as positive. Finally, a population $y(t)$ with unlimited resources grows at a rate proportional to its size, $y'=ry$; if the resources support at most the population $K$, the growth rate per individual decreases as $y$ approaches $K$, and the simplest law of this kind is the \term{logistic equation}
\begin{equation}
y'=ry\Bigl(1-\frac yK\Bigr),\qquad r,K>0 .
\label{eq:logistic}
\end{equation}
All these equations are of the first order, and all except the logistic equation and the equation with quadratic resistance are linear; \cref{sec:separable,sec:linear-first-order} solve them.

\subsection{Two Warnings}

The equation $y'=ky$ has solutions defined for all times, and its initial value problems have exactly one solution. Neither property holds in general, even for very simple right sides.

The equation $y'=y^2$ has a polynomial right side. For $y_0>0$, the function
\[
y(t)=\frac{y_0}{1-y_0t}
\]
satisfies $y'=y_0^2/(1-y_0t)^2=y^2$ and $y(0)=y_0$, but only on the interval $(-\infty,1/y_0)$; as $t$ approaches $1/y_0$ from the left, $y(t)$ tends to $\infty$ (\cref{fig:two-warnings}A). \Cref{sec:separable} shows that there is no other solution, so the solution of this initial value problem \term{blows up} in finite time and cannot be continued beyond $t=1/y_0$. The interval on which the solution exists is not visible in the equation, and it depends on the initial value: the larger $y_0$, the sooner the blow-up.

The equation $y'=\sqrt{|y|}$ with $y(0)=0$ has the solution $y\equiv0$. It has others as well. For every $c\ge0$, the function
\[
y_c(t)\coloneqq\begin{cases}0,&t\le c,\\(t-c)^2/4,&t>c,\end{cases}
\]
is differentiable, with $y_c'(t)=0$ for $t\le c$ and $y_c'(t)=(t-c)/2=\sqrt{y_c(t)}$ for $t>c$, the two one-sided derivatives at $c$ both being $0$. So $y_c$ is a solution on $\mathbb R$ with $y_c(0)=0$, and the initial value problem has infinitely many solutions (\cref{fig:two-warnings}B). A solution may rest at $0$ for an arbitrary time and then leave. The right side is continuous, but its slope with respect to $y$ is unbounded near $y=0$, and \cref{sec:existence-uniqueness} shows that a bound on this slope restores uniqueness.

\begin{figure}
\sdpanel{A}{\begin{tikzpicture}
\begin{axis}[sdplot, width=0.49\textwidth, height=0.42\textwidth, xmin=-1, xmax=2.2, ymin=0, ymax=6,
  xtick={-1,0,0.5,1,2}, xticklabels={$-1$,0,$1/2$,1,2}, ytick={0,2,4,6}, xlabel={$t$}, ylabel={$y$}]
\draw[black!45, densely dotted, line width=0.4pt] (axis cs:0.5,0) -- (axis cs:0.5,6) (axis cs:1,0) -- (axis cs:1,6) (axis cs:2,0) -- (axis cs:2,6);
\addplot[black, line width=0.8pt, domain=-1:0.4167, samples=80] {2/(1-2*x)};
\addplot[black, dashed, line width=0.8pt, domain=-1:0.8333, samples=80] {1/(1-x)};
\addplot[black!45, line width=0.8pt, domain=-1:1.9167, samples=80] {0.5/(1-0.5*x)};
\end{axis}
\end{tikzpicture}}\hfill
\sdpanel{B}{\begin{tikzpicture}
\begin{axis}[sdplot, width=0.49\textwidth, height=0.42\textwidth, xmin=-3, xmax=3, ymin=-1.6, ymax=1.6,
  xtick={-3,-2,-1,0,1,2,3}, ytick={-1,0,1}, xlabel={$t$}, ylabel={$y$}]
\addplot[black, line width=1.1pt, domain=-3:3, samples=2] {0};
\foreach \c in {0,0.5,1}{\addplot[black, line width=0.7pt, domain=\c:3, samples=40] {(x-\c)^2/4};}
\foreach \d in {0,-0.5,-1}{\addplot[black!45, line width=0.7pt, domain=-3:\d, samples=40] {-(\d-x)^2/4};}
\fill[black] (axis cs:0,0) circle (1.4pt);
\end{axis}
\end{tikzpicture}}
\caption{Two warnings. (A) The solutions $y=y_0/(1-y_0t)$ of $y'=y^2$ with $y(0)=y_0$ for $y_0=2$ (solid), $1$ (dashed) and $1/2$ (grey) blow up at $t=1/y_0$ (dotted lines), although the right side is a polynomial. (B) Through the point $(0,0)$ (dot) pass infinitely many solutions of $y'=\sqrt{|y|}$: the zero solution (thick), and solutions that leave the axis at any time $c\ge0$ along the parabolas $(t-c)^2/4$ (black) or arrive on it at any time $d\le0$ along $-(d-t)^2/4$ (grey); any such arrival can be combined with any such departure.\label{fig:two-warnings}}
\end{figure}

The negative solutions can be found in the same way: for $y<0$ the equation reads $y'=\sqrt{-y}$, and $y(t)=-(d-t)^2/4$ for $t<d$ satisfies it, with $y'=(d-t)/2=\sqrt{-y}$. A solution through $(0,0)$ may therefore arrive at the axis from below at any time $d\le0$, rest there, and leave upwards at any time $c\ge0$. Every solution increases, since $y'\ge0$; what is not determined is how long it stays at the value~$0$.

\subsection{The Equivalent Integral Equation}

A differential equation with a continuous right side can be rewritten as an equation that contains no derivatives. The integral form is the starting point of the existence theory of \cref{sec:existence-uniqueness} and of the numerical methods of \cref{sec:euler}.

\begin{proposition}[The integral equation]
\label{prop:integral-equation}
Let $D\subseteq\mathbb R^2$, let $f\colon D\to\mathbb R$ be continuous, let $I$ be an interval containing $t_0$, and let $y\colon I\to\mathbb R$ be a continuous function with $(t,y(t))\in D$ for all $t\in I$. Then $y$ is a solution of the initial value problem $y'=f(t,y)$, $y(t_0)=y_0$, on $I$ if and only if
\begin{equation}
y(t)=y_0+\int_{t_0}^tf\bigl(s,y(s)\bigr)\,ds\qquad\text{for all }t\in I .
\label{eq:integral-equation}
\end{equation}
Every solution of $y'=f(t,y)$ on an interval is continuously differentiable.
\end{proposition}

\begin{proof}
The function $s\mapsto f(s,y(s))$ is continuous on $I$ whenever $y$ is, as a composition of continuous functions (\cref{sec:limits-rn}); so the integral in \eqref{eq:integral-equation} exists for every $t\in I$. Let $y$ be a solution. Its derivative $y'(t)=f(t,y(t))$ is continuous, which proves the last statement, and by the second part of the fundamental theorem of calculus (\cref{thm:ftc2}), applied on the interval between $t_0$ and $t$, $y(t)-y(t_0)=\int_{t_0}^ty'(s)\,ds$, which is \eqref{eq:integral-equation}. Conversely, if $y$ satisfies \eqref{eq:integral-equation}, then $y(t_0)=y_0$, and by the first part of the fundamental theorem (\cref{thm:ftc1}) the right side is differentiable with the derivative $f(t,y(t))$, one-sided at an end of $I$; so $y$ is a solution.
\end{proof}

The integral equation contains the initial value, which the differential equation alone does not, and it makes sense for every continuous $y$, so that one can try to construct a solution as the limit of continuous approximations. The method of successive approximations of \cref{sec:existence-uniqueness} inserts an approximation into the right side of \eqref{eq:integral-equation} to obtain the next one, and Euler's method of \cref{sec:euler} approximates the integral over a short step by a rectangle.

\section{Direction Fields and Integral Curves}
\label{sec:direction-fields}

\subsection{Seeing the Equation}

The equation $y'=f(t,y)$ assigns to every point $(t,y)$ of $D$ a number, $f(t,y)$, which is the slope that the graph of a solution must have if it passes through that point. Drawing a short segment with this slope at many points of the plane makes the equation visible.

\begin{definition}[Direction field, integral curve]
\label{def:direction-field}
The \term{direction field}, or slope field, of the equation $y'=f(t,y)$ assigns to each point $(t,y)\in D$ the line through it with slope $f(t,y)$. The graph $\{(t,y(t)):t\in I\}$ of a solution on an interval $I$ is an \term{integral curve} of the equation. A curve on which $f$ is constant, $f(t,y)=m$, is an \term{isocline}, the curve of equal slopes $m$.
\end{definition}

By definition, an integral curve is a curve that has at each of its points the slope prescribed by the direction field: it is tangent to the field everywhere. One can sketch integral curves by hand by drawing the field and then curves that follow it, much as iron filings show the lines of a magnetic field. Drawing a few isoclines helps, since along each of them all segments are parallel.

\Cref{fig:direction-fields} shows the fields of three equations from \eqref{eq:first-examples}. For $y'=t-y$ (panel~A), the isoclines $t-y=m$ are the lines $y=t-m$ of slope $1$. The isocline with $m=1$, the line $y=t-1$, has the slope $1$ that the field prescribes on it, so it is itself an integral curve, and the field shows that all other integral curves approach it as $t$ increases: above the line the slopes are smaller than $1$, below it larger. \Cref{sec:linear-first-order} confirms this with the formula $y=t-1+Ce^{-t}$. For $y'=y(1-y)$ (panel~B), the right side does not depend on $t$, so the field is the same on every vertical line, and shifting an integral curve horizontally gives another one. The isoclines are horizontal lines; on the lines $y=0$ and $y=1$ the slope is $0$, so these lines are integral curves, and between them the solutions rise from the one towards the other along S-shaped curves, which \cref{sec:separable} computes. For $y'=t^2+y^2$ (panel~C), the isoclines are the circles $t^2+y^2=m$ about the origin, and the slopes are positive away from the origin and grow with the distance from it. Every solution increases, and the field suggests that solutions grow so fast that they blow up, both forwards and backwards in time, as the solutions of $y'=y^2$ did; \cref{sec:existence-uniqueness} proves this for the solution through the origin.

\begin{figure}
\sdpanel{A}{\begin{tikzpicture}
\begin{axis}[sdplot, width=0.47\textwidth, height=0.47\textwidth, xmin=-1, xmax=3, ymin=-2, ymax=2,
  xtick={-1,0,1,2,3}, ytick={-2,-1,0,1,2}, xlabel={$t$}, ylabel={$y$}]
\pgfplotsforeachungrouped \tt in {-0.8,-0.4,...,2.85}{%
 \pgfplotsforeachungrouped \yy in {-1.8,-1.4,...,1.85}{%
  \pgfmathsetmacro\m{\tt-\yy}\pgfmathsetmacro\dx{0.12/sqrt(1+\m*\m)}\pgfmathsetmacro\dy{\m*\dx}%
  \edef\temp{\noexpand\draw[black!45, line width=0.4pt] (axis cs:\tt-\dx,\yy-\dy) -- (axis cs:\tt+\dx,\yy+\dy);}\temp}}
\foreach \c in {-2,-1,-0.4,0.4,1,2}{\addplot[black, line width=0.7pt, domain=-1:3, samples=60] {x-1+\c*exp(-x)};}
\addplot[black, line width=0.7pt, dashed, domain=-1:3, samples=2] {x-1};
\end{axis}
\end{tikzpicture}}\hfill
\sdpanel{B}{\begin{tikzpicture}
\begin{axis}[sdplot, width=0.47\textwidth, height=0.47\textwidth, xmin=0, xmax=6, ymin=-0.5, ymax=1.5,
  xtick={0,2,4,6}, ytick={-0.5,0,0.5,1,1.5}, yticklabels={$-1/2$,0,$1/2$,1,$3/2$}, xlabel={$t$}, ylabel={$y$}]
\pgfplotsforeachungrouped \tt in {0.3,0.9,...,5.75}{%
 \pgfplotsforeachungrouped \yy in {-0.4,-0.2,...,1.45}{%
  \pgfmathsetmacro\m{\yy*(1-\yy)}\pgfmathsetmacro\dx{0.18/sqrt(1+9*\m*\m)}\pgfmathsetmacro\dy{\m*\dx}%
  \edef\temp{\noexpand\draw[black!45, line width=0.4pt] (axis cs:\tt-\dx,\yy-\dy) -- (axis cs:\tt+\dx,\yy+\dy);}\temp}}
\addplot[black, line width=0.7pt, dashed, domain=0:6, samples=2] {0};
\addplot[black, line width=0.7pt, dashed, domain=0:6, samples=2] {1};
\foreach \c in {0.03,0.15,0.5}{\addplot[black, line width=0.7pt, domain=0:6, samples=60] {1/(1+(1-\c)/\c*exp(-x))};}
\addplot[black, line width=0.7pt, domain=0:6, samples=60] {1/(1-(1/3)*exp(-x))};
\addplot[black, line width=0.7pt, domain=0:0.6931, samples=40] {1/(1-6*exp(-x))};
\addplot[black, line width=0.7pt, domain=0:1.9459, samples=60] {1/(1-21*exp(-x))};
\end{axis}
\end{tikzpicture}}\hspace*{\fill}
\sdpanelrow
\sdfloat{\begin{tabular}[b]{@{}l@{}}{\small\bfseries C}\\[-1pt]\begin{tikzpicture}
\begin{axis}[sdplot, width=0.47\textwidth, height=0.47\textwidth, xmin=-2, xmax=2, ymin=-2, ymax=2,
  xtick={-2,-1,0,1,2}, ytick={-2,-1,0,1,2}, xlabel={$t$}, ylabel={$y$}]
\pgfplotsforeachungrouped \tt in {-1.8,-1.4,...,1.85}{%
 \pgfplotsforeachungrouped \yy in {-1.8,-1.4,...,1.85}{%
  \pgfmathsetmacro\m{\tt*\tt+\yy*\yy}\pgfmathsetmacro\dx{0.12/sqrt(1+\m*\m)}\pgfmathsetmacro\dy{\m*\dx}%
  \edef\temp{\noexpand\draw[black!45, line width=0.4pt] (axis cs:\tt-\dx,\yy-\dy) -- (axis cs:\tt+\dx,\yy+\dy);}\temp}}
\addplot[black, densely dotted, line width=0.6pt, domain=0:360, samples=60] ({cos(x)},{sin(x)});
\addplot[black, line width=0.7pt, smooth] coordinates {(-0.23,-2.301) (-0.16,-1.975) (-0.08,-1.705) (0.00,-1.500) (0.08,-1.339) (0.16,-1.208) (0.24,-1.099) (0.32,-1.004) (0.40,-0.920) (0.48,-0.842) (0.56,-0.769) (0.64,-0.697) (0.72,-0.625) (0.80,-0.551) (0.88,-0.473) (0.96,-0.391) (1.04,-0.301) (1.12,-0.202) (1.20,-0.093) (1.28,0.030) (1.36,0.171) (1.44,0.333) (1.52,0.523) (1.60,0.750) (1.68,1.028) (1.76,1.380) (1.84,1.846) (1.90,2.300)};
\addplot[black, line width=0.7pt, smooth] coordinates {(-0.75,-2.301) (-0.72,-2.125) (-0.64,-1.784) (-0.56,-1.536) (-0.48,-1.349) (-0.40,-1.203) (-0.32,-1.088) (-0.24,-0.995) (-0.16,-0.919) (-0.08,-0.855) (0.00,-0.800) (0.08,-0.752) (0.16,-0.708) (0.24,-0.667) (0.32,-0.627) (0.40,-0.587) (0.48,-0.546) (0.56,-0.502) (0.64,-0.455) (0.72,-0.403) (0.80,-0.346) (0.88,-0.281) (0.96,-0.209) (1.04,-0.126) (1.12,-0.032) (1.20,0.075) (1.28,0.200) (1.36,0.345) (1.44,0.517) (1.52,0.723) (1.60,0.976) (1.68,1.293) (1.76,1.709) (1.84,2.284) (1.84,2.301)};
\addplot[black, line width=0.7pt, smooth] coordinates {(-1.63,-2.302) (-1.60,-2.076) (-1.52,-1.612) (-1.44,-1.272) (-1.36,-1.011) (-1.28,-0.806) (-1.20,-0.641) (-1.12,-0.507) (-1.04,-0.397) (-0.96,-0.307) (-0.88,-0.234) (-0.80,-0.174) (-0.72,-0.126) (-0.64,-0.088) (-0.56,-0.059) (-0.48,-0.037) (-0.40,-0.021) (-0.32,-0.011) (-0.24,-0.005) (-0.16,-0.001) (-0.08,-0.000) (0.00,0.000) (0.08,0.000) (0.16,0.001) (0.24,0.005) (0.32,0.011) (0.40,0.021) (0.48,0.037) (0.56,0.059) (0.64,0.088) (0.72,0.126) (0.80,0.174) (0.88,0.234) (0.96,0.307) (1.04,0.397) (1.12,0.507) (1.20,0.641) (1.28,0.806) (1.36,1.011) (1.44,1.272) (1.52,1.612) (1.60,2.076) (1.63,2.302)};
\addplot[black, line width=0.7pt, smooth] coordinates {(-1.81,-2.304) (-1.76,-1.895) (-1.68,-1.438) (-1.60,-1.095) (-1.52,-0.827) (-1.44,-0.610) (-1.36,-0.432) (-1.28,-0.282) (-1.20,-0.155) (-1.12,-0.046) (-1.04,0.047) (-0.96,0.128) (-0.88,0.198) (-0.80,0.258) (-0.72,0.311) (-0.64,0.357) (-0.56,0.397) (-0.48,0.433) (-0.40,0.465) (-0.32,0.493) (-0.24,0.520) (-0.16,0.546) (-0.08,0.572) (0.00,0.600) (0.08,0.630) (0.16,0.665) (0.24,0.706) (0.32,0.755) (0.40,0.814) (0.48,0.888) (0.56,0.979) (0.64,1.093) (0.72,1.239) (0.80,1.426) (0.88,1.674) (0.96,2.012) (1.01,2.300)};
\end{axis}
\end{tikzpicture}\end{tabular}}{Direction fields (grey segments) and integral curves (black). (A) $y'=t-y$; the dashed line $y=t-1$ is an isocline that is also an integral curve, and all other integral curves approach it. (B) $y'=y(1-y)$; the field does not depend on $t$, the lines $y=0$ and $y=1$ (dashed) are integral curves, and the solutions between them rise from one to the other. (C) $y'=t^2+y^2$; the isocline of slope $1$ is the unit circle (dotted), and the integral curves, through $(0,c)$ for $c=-1.5$, $-0.8$, $0$ and $0.6$, computed numerically, leave the window steeply in both directions.\label{fig:direction-fields}}
\end{figure}

\subsection{Integral Curves Do Not Cross}

Two integral curves of the same equation may meet at a point. At that point both have the slope $f(t,y)$, so they are tangent to each other there; in \cref{fig:two-warnings}B the parabolas touch the axis in this way. But can they meet at all? If the initial value problem with the initial point $(t_1,y_1)$ has only one solution, two solutions through $(t_1,y_1)$ are restrictions of the same function, and their integral curves coincide on their common interval. \Cref{sec:existence-uniqueness} proves that this is the case whenever $\partial f/\partial y$ is continuous on an open set $D$ (\cref{cor:local-uniqueness}). Then two integral curves that meet coincide over their common interval, and the integral curves fill the region $D$ like the fibres of a fabric, one through each point.

This has a useful consequence for estimates. If $y_1$ and $y_2$ are solutions with $y_1(t_0)<y_2(t_0)$, then $y_1(t)<y_2(t)$ wherever both are defined, since by the intermediate value theorem (\cref{cor:ivt}) a change of order would require a point where they are equal. A solution that starts between two others is trapped between them. In \cref{fig:direction-fields}B, a solution with $0<y(t_0)<1$ stays between the constant solutions $0$ and $1$ for all time, and a solution with $y(t_0)>1$ stays above $1$; this is how the qualitative theory of \cref{sec:autonomous} argues.

\begin{note}{Equations that cannot be solved by quadratures}
The equation $y'=t^2+y^2$ of \cref{fig:direction-fields}C is a \emph{Riccati equation}, a first-order equation whose right side is a quadratic polynomial in $y$ with coefficients depending on $t$. Liouville proved in 1841 that the equation $y'=y^2+t^m$ can be solved by quadratures, that is, by finitely many algebraic operations, integrations and exponentials applied to elementary functions, only when $m=-2$ or $m=-4k/(2k\pm1)$ with $k\in\mathbb N$~\cite{liouville1841}; the value $k=0$ gives the separable case $m=0$, and as $k\to\infty$ the exceptional values tend to $-2$. Watson gives an account of the result and of its history~\cite[Chap.~IV]{watson1944}. All exceptional values are at most $0$, so $m=2$ is not among them, the solutions of $y'=t^2+y^2$ are not elementary. The theory is that of \cref{thm:liouville}, the differential algebra of Liouville, Ritt and Rosenlicht. The substitution $y=-w'/w$, which \cref{sec:substitutions} explains, turns the equation into the linear second-order equation $w''+t^2w=0$, whose solutions are expressed by Bessel functions~\cite[Chap.~IV]{watson1944}; equations of this kind are solved by power series in Chapter~26. The failure of formulas is the rule, not the exception: almost every equation met in applications has to be studied with the qualitative and numerical methods of the later sections.
\end{note}

\section{Separable Equations}
\label{sec:separable}

\subsection{Separation of Variables}

An equation of the form
\begin{equation}
y'=g(t)\,h(y),
\label{eq:separable}
\end{equation}
in which the right side is a product of a function of $t$ and a function of $y$, is called \term{separable}. The equations $y'=ky$, $y'=y(1-y)$, $y'=y^2$ and $y'=\sqrt{|y|}$ of \cref{sec:ode-basics} are separable, with $g\equiv1$; so are the laws of cooling and of falling, and every autonomous equation $y'=f(y)$. Leibniz's notation suggests a recipe. Write $y'$ as $dy/dt$, move all terms with $y$ to one side and all terms with $t$ to the other, $dy/h(y)=g(t)\,dt$, and integrate both sides:
\[
\int\frac{dy}{h(y)}=\int g(t)\,dt+C .
\]
Solving the resulting equation for $y$ gives the solutions. The manipulation of $dy$ and $dt$ as if they were numbers needs a justification, and the theorem below provides it; it also shows where the recipe can go wrong, namely at the zeros of $h$, by which the recipe divides.

\begin{theorem}[Separation of variables]
\label{thm:separable}
Let $J$ and $K$ be open intervals, let $g\colon J\to\mathbb R$ and $h\colon K\to\mathbb R$ be continuous, and let $t_0\in J$ and $y_0\in K$ with $h(y_0)\ne0$. Let $K_0\subseteq K$ be the largest open interval containing $y_0$ on which $h$ has no zero, and let
\[
G(t)\coloneqq\int_{t_0}^tg(s)\,ds\quad(t\in J),\qquad H(y)\coloneqq\int_{y_0}^y\frac{d\eta}{h(\eta)}\quad(y\in K_0).
\]
\begin{enumerate}
\item $H$ is continuously differentiable and strictly monotone, and it maps $K_0$ onto an open interval $H(K_0)$ that contains $0$.
\item There is an open interval $I\subseteq J$ containing $t_0$ with $G(I)\subseteq H(K_0)$, and on every such interval the function $y(t)\coloneqq H^{-1}(G(t))$ is a solution of the initial value problem $y'=g(t)h(y)$, $y(t_0)=y_0$.
\item Every solution of this initial value problem on an interval containing $t_0$ whose values lie in $K_0$ is of the form $H^{-1}\circ G$; that is, it satisfies $H(y(t))=G(t)$.
\item If $c\in K$ and $h(c)=0$, the constant function $y\equiv c$ is a solution on $J$.
\end{enumerate}
\end{theorem}

\begin{proof}
(i) The function $h$ is continuous and has no zero on the interval $K_0$, so it has a constant sign there, by the intermediate value theorem (\cref{cor:ivt}), and $1/h$ is continuous on $K_0$. By \cref{thm:ftc1}, $H$ is differentiable with $H'=1/h$, a continuous function of constant sign, so $H$ is strictly monotone by \cref{cor:monotone-derivative}. The image of the interval $K_0$ under the continuous function $H$ is an interval (\cref{cor:interval-image}); it is open, because a strictly monotone function on an open interval attains no largest and no smallest value, and it contains $H(y_0)=0$.

(ii) Since $G$ is continuous with $G(t_0)=0$ and $H(K_0)$ is an open interval containing $0$, there is $\delta>0$ such that $G(t)\in H(K_0)$ for $|t-t_0|<\delta$; this gives an interval $I$. On any such interval, $y=H^{-1}\circ G$ is defined. The inverse $H^{-1}\colon H(K_0)\to K_0$ is differentiable, with $(H^{-1})'(u)=1/H'(H^{-1}(u))=h(H^{-1}(u))$, by \cref{thm:inverse-derivative}, since $H'$ never vanishes; and $G'=g$ by \cref{thm:ftc1}. By the chain rule,
\[
y'(t)=(H^{-1})'\bigl(G(t)\bigr)\,G'(t)=h\bigl(y(t)\bigr)\,g(t),
\]
and $y(t_0)=H^{-1}(0)=y_0$.

(iii) Let $y$ be a solution on an interval $I\ni t_0$ with values in $K_0$. By the chain rule, $H\circ y$ is differentiable, with
\[
\frac d{dt}H\bigl(y(t)\bigr)=H'\bigl(y(t)\bigr)\,y'(t)=\frac{g(t)\,h(y(t))}{h(y(t))}=g(t)=G'(t).
\]
So $H\circ y-G$ has derivative zero on $I$ and is constant (\cref{cor:zero-derivative}), with the value $H(y_0)-G(t_0)=0$ at $t_0$. Hence $H(y(t))=G(t)$ and $y(t)=H^{-1}(G(t))$.

(iv) The derivative of a constant is $0=g(t)h(c)$.
\end{proof}

The theorem gives a solution on the interval of all $t$ near $t_0$ for which $G(t)$ lies in $H(K_0)$, and it says that no other solution exists as long as the values stay away from the zeros of $h$. In practice one computes the two antiderivatives, writes down the equation $H(y)=G(t)$, solves it for $y$, and then checks for which $t$ the result is defined. The zeros of $h$ need separate attention. Each zero $c$ gives a constant solution, part (iv), which the recipe loses when it divides by $h(y)$. Whether a nonconstant solution can reach the value $c$ in finite time is not decided by the theorem: for $y'=\sqrt{|y|}$ the solutions $y_c$ of \cref{sec:ode-basics} do reach $0$, while for $y'=y^2$ and $y'=ky$ no solution reaches $0$ unless it is $0$ throughout. \Cref{sec:existence-uniqueness} shows that the second behaviour is the rule when $h$ is continuously differentiable, and \cref{sec:autonomous} computes when a solution reaches a zero of $h$ in finite time.

For the equation $y'=y^2$ with $y(0)=y_0>0$, the theorem applies with $g\equiv1$, $h(y)=y^2$ and $K_0=(0,\infty)$:
\[
H(y)=\int_{y_0}^y\frac{d\eta}{\eta^2}=\frac1{y_0}-\frac1y,\qquad G(t)=t ,
\]
and $H$ maps $(0,\infty)$ onto $(-\infty,1/y_0)$. So the solution is defined exactly for $t<1/y_0$, where $1/y=1/y_0-t$, and it is the function of \cref{fig:two-warnings}A; by (iii) every solution with values in $(0,\infty)$ agrees with it, and \cref{cor:local-uniqueness} below excludes solutions that reach $0$. The recipe in this example reads $dy/y^2=dt$, $-1/y=t+C$, $y=-1/(t+C)$, and it loses the constant solution $y\equiv0$.

\subsection{Exponential Growth and Decay}

For $y'=ky$ with $y_0>0$ the theorem gives $H(y)=\ln(y/y_0)$ and $G(t)=k(t-t_0)$, hence $y=y_0e^{k(t-t_0)}$, in agreement with \cref{thm:exp-ode}. \Cref{sec:exp-log-numerics} drew the consequences: for $k>0$ the solution doubles in every interval of length $\ln2/k$, and for $k=-\lambda<0$ it halves in every interval of length $T_{1/2}=\ln2/\lambda$, the half-life, so that the fraction $2^{-t/T_{1/2}}$ remains after the time $t$.

Radiocarbon dating rests on this law. It assumes that a living organism contains the radioactive isotope carbon-14 in the same small proportion as the carbon of the atmosphere, and that after its death this proportion decays with the half-life of $5730$ years~\cite{godwin1962}. If a sample contains the fraction $p$ of the carbon-14 proportion of living matter, its age $t$ satisfies $2^{-t/T_{1/2}}=p$, that is,
\[
t=\frac{T_{1/2}}{\ln2}\,\ln\frac1p .
\]
For $p=0.30$ this gives about $9950$ years, and for $p=1/4$ exactly two half-lives, $11\,460$ years. The formula also says how accurate the age is. Its derivative with respect to $p$ is $-T_{1/2}/(p\ln2)$, so a relative error $\delta$ in the measured fraction $p$ produces an absolute error of about $T_{1/2}\delta/\ln2$ in the age, independently of $p$: a measurement accurate to one per cent determines the age to within about $83$ years, whether the sample is old or young. The reason is that the logarithm converts relative changes into absolute ones: $\ln(p(1+\delta))-\ln p=\ln(1+\delta)\approx\delta$.

Newton's law of cooling $T'=-k(T-A)$ becomes the decay law for the excess temperature $u\coloneqq T-A$, which satisfies $u'=-ku$. So $T(t)=A+(T_0-A)e^{-kt}$: the excess temperature decays exponentially. A cup of coffee at $90$ degrees in a room at $20$ degrees that has cooled to $60$ degrees after $10$ minutes has $e^{-10k}=40/70$, so $k=\ln(7/4)/10=0.0560$ per minute, and it reaches $40$ degrees when $e^{-kt}=20/70$, after $t=10\ln(7/2)/\ln(7/4)=22.4$ minutes.

\subsection{The Logistic Equation}

The logistic equation \eqref{eq:logistic} is separable with $g\equiv r$ and $h(y)=y(1-y/K)$, which vanishes at $0$ and $K$. These are the constant solutions. For $0<y_0<K$ we have $K_0=(0,K)$, and the partial fraction decomposition (\cref{thm:partial-fractions})
\[
\frac1{y(1-y/K)}=\frac K{y(K-y)}=\frac1y+\frac1{K-y}
\]
gives, for $0<y<K$,
\[
H(y)=\int_{y_0}^y\Bigl(\frac1\eta+\frac1{K-\eta}\Bigr)d\eta=\ln\frac y{K-y}-\ln\frac{y_0}{K-y_0} .
\]
$H$ maps $(0,K)$ onto all of $\mathbb R$, so the solution exists for all $t$, and $H(y)=r(t-t_0)$ is solved by
\begin{equation}
y(t)=\frac K{1+Ce^{-r(t-t_0)}},\qquad C\coloneqq\frac{K-y_0}{y_0}.
\label{eq:logistic-solution}
\end{equation}
The same formula, checked by differentiation, gives the solutions for the other initial values $y_0\ne0$, on the interval of those $t$ near $t_0$ for which the denominator does not vanish. For $y_0>K$ the constant $C$ lies in $(-1,0)$, the denominator is positive for $t>t_0+\ln(-C)/r$, and the solution decreases towards $K$; going backwards it blows up at the time $t_0+\ln(-C)/r$, which is earlier than $t_0$. For $y_0<0$, $C<-1$, and the solution decreases and tends to $-\infty$ at the time $t_0+\ln(-C)/r>t_0$; negative populations have no meaning, but the formula is still the solution of the equation.

\Cref{fig:logistic} shows the solutions. Those that start between $0$ and $K$ trace an S-shaped curve: almost exponential growth at first, while $y$ is small and $y'\approx ry$, then a slowing down, and finally saturation at the \term{carrying capacity} $K$. The curve is symmetric about the point where $y=K/2$, and this is its point of inflection: differentiating the equation gives $y''=r(1-2y/K)\,y'$, which changes sign exactly when $y$ passes $K/2$. Writing $C=e^{r(t_m-t_0)}$, the solution is $y(t)=K\sigma(r(t-t_m))$ with the logistic function $\sigma(x)=1/(1+e^{-x})$ of \cref{sec:hyperbolic}; the identity $\sigma'=\sigma(1-\sigma)$ noted there is the logistic equation with $r=K=1$, and the inflection time $t_m$ is the time at which the population reaches half its capacity.

\begin{figure}
\sdfloat{%
\begin{tikzpicture}
\begin{axis}[sdplot, width=0.6\textwidth, height=0.42\textwidth, xmin=0, xmax=9, ymin=0, ymax=2.1,
  xtick={0,2,4,6,8}, ytick={0,0.5,1,1.5,2}, yticklabels={0,$K/2$,$K$,$3K/2$,$2K$}, xlabel={$t$}, ylabel={$y$}]
\addplot[black, dashed, line width=0.6pt, domain=0:9, samples=2] {1};
\addplot[black!45, densely dotted, line width=0.5pt, domain=0:9, samples=2] {0.5};
\foreach \c in {0.02,0.1,0.3,0.7}{\addplot[black, line width=0.8pt, domain=0:9, samples=80] {1/(1+(1-\c)/\c*exp(-x))};}
\foreach \c in {1.5,2.0}{\addplot[black!45, line width=0.8pt, domain=0:9, samples=80] {1/(1+(1-\c)/\c*exp(-x))};}
\end{axis}
\end{tikzpicture}}{Solutions of the logistic equation $y'=ry(1-y/K)$ with $r=1$, for $y(0)=0.02K$, $0.1K$, $0.3K$, $0.7K$ (black) and $1.5K$, $2K$ (grey). Solutions below the carrying capacity $K$ (dashed) rise along S-shaped curves with their point of inflection at $y=K/2$ (dotted); those above it decrease towards $K$.\label{fig:logistic}}
\end{figure}

The logistic law was proposed for human populations by Verhulst in 1838~\cite{verhulst1838}. Its two constants have a simple meaning: $r$ is the growth rate of a small population, and $K$ the population that the environment sustains. Together with the constant $C$, they are determined by three observations at equally spaced times $t_0$, $t_0+h$ and $t_0+2h$: by \eqref{eq:logistic-solution}, $1/y=1/K+(C/K)e^{-r(t-t_0)}$, so the two differences of the observed values of $1/y$ have the ratio $e^{-rh}$, which gives $r$, and then $K$ and $C$ follow from two linear equations.

\subsection{Leaking Tanks and Falling Bodies}

A cylindrical tank with cross-section $A$ drains through a small hole of area $a$ in its bottom. By Torricelli's law, water leaves through the hole with the speed $\sqrt{2gH}$ when its depth is $H$, so the volume $AH$ decreases at the rate $a\sqrt{2gH}$, and
\[
H'=-k\sqrt H,\qquad k\coloneqq\frac aA\sqrt{2g} .
\]
For $H_0>0$ separation of variables gives $2(\sqrt H-\sqrt{H_0})=-kt$, that is,
\[
H(t)=\Bigl(\sqrt{H_0}-\frac{kt}2\Bigr)^2\qquad\text{for }0\le t\le T_{\mathrm e}\coloneqq\frac{2\sqrt{H_0}}k ,
\]
and at the time $T_{\mathrm e}$ the tank is empty. Unlike the solutions of the decay law, this solution reaches the zero of the right side in finite time, and it continues as $H\equiv0$ afterwards; the glued function is differentiable at $T_{\mathrm e}$, where both one-sided derivatives vanish. The equation is the mirror image of $y'=\sqrt{|y|}$, and it shows the non-uniqueness of \cref{fig:two-warnings}B as a physical phenomenon. Run backwards in time, the initial value problem with $H(t_1)=0$ has the solutions that are $0$ for $t\ge c$ and $(k/2)^2(c-t)^2$ for $t<c$, one for each $c\le t_1$. An empty tank does not reveal when it became empty: the present state does not determine the past.

A body falling in air with the resistance $kv^2$ satisfies $v'=g-cv^2$ with $c\coloneqq k/m$. The right side vanishes at the \term{terminal velocity} $v_{\mathrm T}\coloneqq\sqrt{g/c}$, a constant solution, and for $|v|<v_{\mathrm T}$ we have $g-cv^2=c(v_{\mathrm T}^2-v^2)$. By the derivative of the inverse hyperbolic tangent (\cref{prop:inverse-hyperbolic}),
\[
\frac d{dv}\,\frac1{cv_{\mathrm T}}\operatorname{artanh}\frac v{v_{\mathrm T}}=\frac1{c(v_{\mathrm T}^2-v^2)} ,
\]
so for a body released at rest, $\operatorname{artanh}(v/v_{\mathrm T})=cv_{\mathrm T}t$, and with $cv_{\mathrm T}=g/v_{\mathrm T}$,
\[
v(t)=v_{\mathrm T}\tanh\frac{gt}{v_{\mathrm T}} .
\]
The speed rises like $gt$ at first, as in free fall, and approaches $v_{\mathrm T}$ exponentially, since $1-\tanh x=2/(e^{2x}+1)<2e^{-2x}$. It reaches $99$ per cent of the terminal velocity when $gt/v_{\mathrm T}=\operatorname{artanh}0.99=2.647$; for a terminal velocity of $50$ metres per second and $g=9.81$ metres per second squared, this takes $13.5$ seconds. Integrating once more gives the distance fallen, $(v_{\mathrm T}^2/g)\ln\cosh(gt/v_{\mathrm T})$.

\subsection{The Catenary}
\label{subsec:catenary}

\Cref{sec:hyperbolic} stated that a hanging chain takes the shape of the curve $y=a\cosh(x/a)$, and that this follows from the equation $y''=\sqrt{1+y'^2}/a$. We now derive the equation and solve it. Here a solution is a curve in the plane, so we use the coordinates $x$ and $y$, with $y'=dy/dx$.

A flexible chain of weight $w$ per unit length hangs between two points; put its lowest point on the $y$-axis. Cut out the piece of chain between the lowest point and the point $(x,y(x))$ with $x>0$, and consider the forces on it (\cref{fig:catenary}). A flexible chain transmits forces only along itself, so at the lowest point the rest of the chain pulls the piece horizontally, with a force $H>0$, and at the point $(x,y(x))$ it pulls along the tangent, with a force of some magnitude $T$ and direction $(\cos\theta,\sin\theta)$, where $\tan\theta=y'(x)$. The third force is the weight $ws(x)$ of the piece, where $s(x)=\int_0^x\sqrt{1+y'(\xi)^2}\,d\xi$ is its length (\cref{thm:arc-length}). The piece is at rest, so the forces balance: horizontally $T\cos\theta=H$ and vertically $T\sin\theta=ws(x)$. Dividing,
\[
y'(x)=\tan\theta=\frac{s(x)}a,\qquad a\coloneqq\frac Hw ,
\]
and differentiating with $s'(x)=\sqrt{1+y'(x)^2}$ gives the equation of the catenary,
\begin{equation}
y''=\frac1a\sqrt{1+y'^2} .
\label{eq:catenary}
\end{equation}
It is of the second order, but $y$ itself does not occur in it, so it is a first-order equation for the slope $p\coloneqq y'$:
\[
p'=\frac1a\sqrt{1+p^2},\qquad p(0)=0 ,
\]
since the tangent is horizontal at the lowest point. This equation is separable with $g\equiv1/a$ and $h(p)=\sqrt{1+p^2}$, which has no zero, so $K_0=\mathbb R$. The antiderivative of $1/\sqrt{1+p^2}$ is $\operatorname{arsinh}p$, by \cref{prop:inverse-hyperbolic}, so $H(p)=\operatorname{arsinh}p$, which maps $\mathbb R$ onto $\mathbb R$, and $G(x)=x/a$. By \cref{thm:separable}, the unique solution is defined on all of $\mathbb R$ and satisfies $\operatorname{arsinh}p=x/a$, that is,
\[
y'(x)=p(x)=\sinh\frac xa,\qquad y(x)=a\cosh\frac xa+C .
\]
Choosing the vertical coordinate so that the lowest point is $(0,a)$ gives $C=0$ and the catenary $y=a\cosh(x/a)$. The derivation also yields the length of chain from the lowest point to $x$, $s(x)=ay'(x)=a\sinh(x/a)$, and the tension $T=H/\cos\theta=H\sqrt{1+y'^2}=H\cosh(x/a)=wy$, which is proportional to the height above the level $y=0$. The constant $a=H/w$ is determined by the length of the chain and the positions of its ends.

\begin{figure}
\sdfloat{%
\begin{tikzpicture}
\begin{axis}[sdplot, width=0.58\textwidth, height=0.55\textwidth, xmin=-2.1, xmax=2.1, ymin=0, ymax=4,
  xtick={-2,-1,0,1,2}, ytick={0,1,2,3,4}, xlabel={$x$}, ylabel={$y$}]
\addplot[black!25, line width=3.5pt, domain=0:1.5, samples=40] {cosh(x)};
\addplot[black, line width=0.9pt, domain=-2.1:2.1, samples=80] {cosh(x)};
\addplot[black, dashed, line width=0.7pt, domain=-2.1:2.1, samples=40] {1+x^2/2};
\draw[black, line width=0.6pt, -{Latex[length=1.8mm]}] (axis cs:0,1) -- (axis cs:-0.9,1) node[anchor=south, font=\scriptsize] {$H$};
\draw[black, line width=0.6pt, -{Latex[length=1.8mm]}] (axis cs:1.5,2.3524) -- (axis cs:1.925,3.257) node[anchor=south east, font=\scriptsize] {$T$};
\draw[black, line width=0.6pt, -{Latex[length=1.8mm]}] (axis cs:0.75,1.2947) -- (axis cs:0.75,0.4) node[anchor=west, font=\scriptsize] {$ws(x)$};
\fill[black] (axis cs:0,1) circle (1.3pt) (axis cs:1.5,2.3524) circle (1.3pt);
\end{axis}
\end{tikzpicture}}{The catenary $y=a\cosh(x/a)$ with $a=1$ (solid) and the parabola $1+x^2/2$ with the same lowest point and curvature (dashed). The shaded piece of chain, from the lowest point to $x=1.5$, is held by the horizontal pull $H$ at its lowest point and the pull $T$ along the tangent at its upper end, and carries its weight $ws(x)$; the balance of these forces gives the equation $y''=\sqrt{1+y'^2}/a$.\label{fig:catenary}}
\end{figure}

\section{Linear Equations}
\label{sec:linear-first-order}

\subsection{The Homogeneous Equation}

A \term{linear} equation of the first order has the form $a(t)y'+b(t)y=c(t)$. Where $a(t)\ne0$ we may divide by it, and we shall always assume that this has been done:
\begin{equation}
y'+p(t)\,y=q(t),
\label{eq:linear}
\end{equation}
with functions $p$ and $q$ that are continuous on an interval $I$. The equation is \term{homogeneous} if $q\equiv0$ and \term{inhomogeneous} otherwise; $q$ is often called the \term{input}, or forcing term. At points where the original coefficient $a$ vanishes, the equation is singular, and its solutions may behave quite differently there: the solutions of $ty'+y=0$ for $t\ne0$ are the functions $c/t$, which have no limit at $0$ unless $c=0$.

\begin{theorem}[The homogeneous linear equation]
\label{thm:linear-homogeneous}
Let $p$ be continuous on an interval $I$, let $t_0\in I$ and $P(t)\coloneqq\int_{t_0}^tp(s)\,ds$. The solutions of $y'+p(t)y=0$ on $I$ are exactly the functions
\[
y(t)=c\,e^{-P(t)},\qquad c\in\mathbb R .
\]
They form a subspace of dimension $1$ of $C^1(I)$, the kernel of the linear map $L\colon C^1(I)\to C(I)$, $Ly\coloneqq y'+py$. The initial value problem with $y(t_0)=y_0$ has exactly one solution, $y_0e^{-P(t)}$, and a solution that vanishes at one point of $I$ vanishes everywhere.
\end{theorem}

\begin{proof}
By \cref{thm:ftc1}, $P$ is differentiable with $P'=p$, so $ce^{-P}$ has the derivative $-pce^{-P}$ by the chain rule and is a solution. Conversely, let $y$ be a solution and $g(t)\coloneqq y(t)e^{P(t)}$. By the product rule,
\[
g'(t)=\bigl(y'(t)+p(t)y(t)\bigr)e^{P(t)}=0 ,
\]
so $g$ is constant on $I$ (\cref{cor:zero-derivative}), $g\equiv g(t_0)=y(t_0)$, and $y=y(t_0)e^{-P}$. The map $L$ is linear by the rules of differentiation, so its kernel is a subspace (\cref{prop:kernel-image}); we have just shown that it is spanned by the function $e^{-P}$, which is not zero, so its dimension is $1$. The remaining statements follow from the formula, since $e^{-P}$ has no zero.
\end{proof}

For a constant $p\equiv-k$ this is \cref{thm:exp-ode}, with the same proof: multiplying by a suitable exponential produces a function with derivative zero. In the language of \cref{ch:linear-maps}, the theorem determines the kernel of the differential operator $L=D+p$, where $D$ is differentiation; the eigenfunctions of $D$ in \cref{sec:eigen-definitions} were the case of a constant $p$. That the kernel has dimension $1$ for an equation of the first order is the first instance of a general fact, which \cref{ch:linear-higher-order} proves: the solutions of a homogeneous linear equation of order $k$ form a space of dimension $k$.

\subsection{The Integrating Factor}

The inhomogeneous equation is solved by the same multiplication. With $\mu(t)\coloneqq e^{P(t)}$, which satisfies $\mu'=p\mu$, the product rule gives, for every differentiable $y$,
\begin{equation}
(\mu y)'=\mu y'+p\mu y=\mu\,(y'+py).
\label{eq:integrating-factor}
\end{equation}
Multiplication by $\mu$ turns the left side of the equation into an exact derivative, and $\mu$ is therefore called an \term{integrating factor}.

\begin{theorem}[Linear equations of the first order]
\label{thm:linear-first-order}
Let $p$ and $q$ be continuous on an interval $I$, let $t_0\in I$, $y_0\in\mathbb R$ and $P(t)\coloneqq\int_{t_0}^tp(s)\,ds$. The initial value problem $y'+p(t)y=q(t)$, $y(t_0)=y_0$, has exactly one solution on $I$, namely
\begin{equation}
y(t)=e^{-P(t)}\Bigl(y_0+\int_{t_0}^te^{P(s)}q(s)\,ds\Bigr).
\label{eq:linear-formula}
\end{equation}
The set of all solutions on $I$ is $y_{\mathrm p}+\ker L$, where $y_{\mathrm p}$ is any one solution and $\ker L=\{ce^{-P}:c\in\mathbb R\}$ is the space of solutions of the homogeneous equation.
\end{theorem}

\begin{proof}
Let $\mu\coloneqq e^P$. By \eqref{eq:integrating-factor}, a differentiable function $y$ on $I$ is a solution of \eqref{eq:linear} if and only if $(\mu y)'=\mu q$, since $\mu$ has no zero. The function $\mu q$ is continuous, so by \cref{cor:antiderivative} this holds if and only if $\mu(t)y(t)-\mu(t_0)y(t_0)=\int_{t_0}^t\mu(s)q(s)\,ds$ for all $t\in I$. With $\mu(t_0)=1$ and $y(t_0)=y_0$ this is \eqref{eq:linear-formula}. So the function \eqref{eq:linear-formula} is a solution, and every solution with $y(t_0)=y_0$ equals it. The last statement is \cref{prop:preimage-flat} for the linear map $L$ of \cref{thm:linear-homogeneous}, whose fibre over $q$ is the set of solutions.
\end{proof}

Two features of the theorem deserve emphasis. First, the solution exists on the whole interval on which $p$ and $q$ are continuous: a linear equation cannot blow up, in contrast with $y'=y^2$, and the interval of existence can be read off from the coefficients. Second, the structure \emph{general solution equals particular solution plus general homogeneous solution} is that of linear systems of equations in \cref{thm:solution-structure}, with the same proof, as \cref{ch:linear-systems,ch:linear-maps} anticipated. The equation $y'=t-y$ of \cref{fig:direction-fields}A illustrates it: with $p\equiv1$ the integrating factor is $e^t$, and $(e^ty)'=te^t$ integrates to $e^ty=(t-1)e^t+C$, so $y=t-1+Ce^{-t}$; the particular solution $t-1$ is the integral curve that all others approach, since the homogeneous part $Ce^{-t}$ decays.

The linearity of $L$ gives the \term{principle of superposition}: if $Ly_1=q_1$ and $Ly_2=q_2$, then $L(\alpha y_1+\beta y_2)=\alpha q_1+\beta q_2$. The response to a sum of inputs is the sum of the responses, so a complicated input can be split into simple ones, for instance into sines and cosines, whose responses are computed separately below.

The formula can be found without guessing the integrating factor, by the method of \term{variation of constants}: one seeks a solution of the form $y=c(t)e^{-P(t)}$, replacing the constant of the homogeneous solution by a function. Then $y'+py=c'e^{-P}$, so $y$ is a solution exactly when $c'=e^Pq$, and integrating gives \eqref{eq:linear-formula} again. For a constant coefficient $p\equiv a$ the formula reads
\begin{equation}
y(t)=e^{-a(t-t_0)}y_0+\int_{t_0}^te^{-a(t-s)}q(s)\,ds ,
\label{eq:duhamel}
\end{equation}
which says that the input $q(s)$ at the time $s$ contributes to the solution at the later time $t$ with the weight $e^{-a(t-s)}$, which depends only on the elapsed time $t-s$: for $a>0$ the equation remembers recent inputs and forgets old ones exponentially. The same formula with the matrix exponential, $\mathbf x(t)=e^{tA}\mathbf x_0+\int_0^te^{(t-s)A}\mathbf f(s)\,ds$, solves the inhomogeneous system $\mathbf x'=A\mathbf x+\mathbf f(t)$, as Chapter~25 shows.

Liouville's criterion, from the proof of \cref{thm:liouville}, concerned an equation of this kind: for rational functions $g$ and $h$ with $g$ not constant, $h\,e^g$ has an elementary antiderivative exactly when $R'+g'R=h$ has a rational solution $R$. By \cref{thm:linear-first-order}, all solutions of this linear equation are $R=e^{-g}\bigl(\int he^g+C\bigr)$, which restates that $Re^g$ is an antiderivative of $he^g$; the content of the criterion lies entirely in the word \emph{rational}.

\subsection{Applications}

A tank holds $100$ litres of water in which salt is dissolved. Brine with $0.2$ kilograms of salt per litre flows in at $5$ litres per minute, the mixture is kept uniform by stirring, and it flows out at the same rate. The amount $y(t)$ of salt, in kilograms, changes by the inflow of $5\cdot0.2=1$ kilogram per minute minus the outflow of $5\cdot y/100$ kilograms per minute:
\[
y'=1-\frac y{20} .
\]
With $y(0)=0$, \eqref{eq:duhamel} gives $y(t)=\int_0^te^{-(t-s)/20}\,ds=20\,(1-e^{-t/20})$. The amount approaches $20$ kilograms, the amount at which the tank has the concentration of the inflow, and after one hour it is $20(1-e^{-3})=19.0$ kilograms.

An electric circuit in which a resistor of resistance $R$ and a capacitor of capacitance $C$ are connected in series to a voltage source $E(t)$ obeys $RCv'+v=E(t)$ for the voltage $v$ across the capacitor, and a body falling with linear resistance obeys $mv'+kv=mg$. Both are linear equations with constant coefficients and a time constant, $RC$ and $m/k$, which measures how fast the solution forgets its initial value. With a constant input the solutions are $v=E+(v_0-E)e^{-t/(RC)}$ and $v=mg/k+(v_0-mg/k)e^{-kt/m}$: after one time constant the distance from the equilibrium value has shrunk by the factor $e^{-1}=0.37$, after five time constants by $e^{-5}=0.0067$. The falling body approaches the terminal velocity $mg/k$, at which weight and resistance balance.

\subsection{Periodic Forcing}

Let the input oscillate: for constants $a>0$ and $\omega>0$, consider
\[
y'+ay=\cos\omega t .
\]
A particular solution of the same form as the input can be found by trying $y_{\mathrm p}=\alpha\cos\omega t+\beta\sin\omega t$. Then
\[
y_{\mathrm p}'+ay_{\mathrm p}=(a\alpha+\omega\beta)\cos\omega t+(a\beta-\omega\alpha)\sin\omega t ,
\]
and this equals $\cos\omega t$ for all $t$ when $a\alpha+\omega\beta=1$ and $a\beta-\omega\alpha=0$, that is, $\alpha=a/(a^2+\omega^2)$ and $\beta=\omega/(a^2+\omega^2)$. With the angle $\varphi\in(0,\pi/2)$ with $\tan\varphi=\omega/a$, so that $a=\sqrt{a^2+\omega^2}\cos\varphi$ and $\omega=\sqrt{a^2+\omega^2}\sin\varphi$, the addition formula gives
\[
y_{\mathrm p}(t)=\frac{a\cos\omega t+\omega\sin\omega t}{a^2+\omega^2}=\frac1{\sqrt{a^2+\omega^2}}\cos(\omega t-\varphi),
\]
and by \cref{thm:linear-first-order} every solution is
\[
y(t)=y_{\mathrm p}(t)+\bigl(y_0-y_{\mathrm p}(0)\bigr)e^{-at}
\]
for the initial value $y(0)=y_0$. The second term, the \term{transient}, decays, and every solution approaches the \term{steady state} $y_{\mathrm p}$, which oscillates with the frequency of the input, with the amplitude $1/\sqrt{a^2+\omega^2}$ and the phase lag $\varphi$, a delay of $\varphi/\omega$ in time (\cref{fig:forced}). For $a=1$ and $\omega=2$ the amplitude is $1/\sqrt5=0.447$ and the phase lag $\arctan2=1.107$, a delay of $0.554$. The amplitude decreases as the frequency grows: the equation passes slow oscillations of the input almost unchanged in amplitude, and damps fast ones. An electrical engineer calls such a system a low-pass filter. The computation is shorter with complex numbers: the complex solution $ce^{i\omega t}$ of $y'+ay=e^{i\omega t}$ has $c(i\omega+a)=1$, so $c=1/(a+i\omega)$, and the real part of $e^{i\omega t}/(a+i\omega)$ is $y_{\mathrm p}$, with modulus $1/|a+i\omega|$ and phase lag $\arg(a+i\omega)=\varphi$. For equations of the second order, \cref{ch:linear-higher-order} finds that the amplitude can become very large when the frequency of the input is close to a natural frequency of the system, the phenomenon of resonance; an equation of the first order has no natural frequency, and its amplitude is largest for a constant input.

\begin{figure}
\sdfloat{%
\begin{tikzpicture}
\begin{axis}[sdplot, width=0.66\textwidth, height=0.36\textwidth, xmin=0, xmax=10, ymin=-1.1, ymax=1.6,
  xtick={0,2,4,6,8,10}, ytick={-1,-0.5,0,0.5,1,1.5}, xlabel={$t$}]
\addplot[black!45, densely dotted, line width=0.6pt, domain=0:10, samples=200] {cos(deg(2*x))};
\addplot[black, dashed, line width=0.7pt, domain=0:10, samples=200] {(cos(deg(2*x))+2*sin(deg(2*x)))/5};
\addplot[black, line width=0.8pt, domain=0:10, samples=200] {1.3*exp(-x)+(cos(deg(2*x))+2*sin(deg(2*x)))/5};
\end{axis}
\end{tikzpicture}}{Periodic forcing of $y'+y=\cos2t$. The solution with $y(0)=1.5$ (solid) approaches the steady state $(\cos2t+2\sin2t)/5=0.447\cos(2t-1.107)$ (dashed) as the transient $1.3e^{-t}$ decays. The steady state has the frequency of the input $\cos2t$ (dotted), a smaller amplitude and a delay of $0.554$.\label{fig:forced}}
\end{figure}

\subsection{Dependence on the Data}

The formula \eqref{eq:linear-formula} shows how the solution depends on the initial value and on the input. When the coefficient $p$ is bounded below by a positive constant, the dependence is uniformly stable.

\begin{proposition}[Stability of linear equations]
\label{prop:linear-dependence}
Let $p$, $q$ and $r$ be continuous on $[t_0,\infty)$ with $p(t)\ge\alpha>0$ for all $t$, and let $y$ and $z$ be the solutions of $y'+py=q$, $y(t_0)=y_0$ and $z'+pz=r$, $z(t_0)=z_0$. Then, for all $t\ge t_0$,
\[
|y(t)-z(t)|\le|y_0-z_0|\,e^{-\alpha(t-t_0)}+\frac1\alpha\max_{t_0\le s\le t}|q(s)-r(s)| .
\]
\end{proposition}

\begin{proof}
The difference $w\coloneqq y-z$ solves $w'+pw=q-r$ with $w(t_0)=y_0-z_0$, so by \eqref{eq:linear-formula}
\[
w(t)=e^{-P(t)}w(t_0)+\int_{t_0}^te^{-(P(t)-P(s))}\bigl(q(s)-r(s)\bigr)\,ds .
\]
For $t_0\le s\le t$, $P(t)-P(s)=\int_s^tp\ge\alpha(t-s)$, so $e^{-(P(t)-P(s))}\le e^{-\alpha(t-s)}$, and in particular $e^{-P(t)}\le e^{-\alpha(t-t_0)}$. Bounding $|q-r|$ by its maximum on $[t_0,t]$ and using $\int_{t_0}^te^{-\alpha(t-s)}\,ds=(1-e^{-\alpha(t-t_0)})/\alpha\le1/\alpha$ gives the claim.
\end{proof}

With $q=r$, the proposition says that any two solutions of the same equation approach each other exponentially: the equation forgets its initial value, which justifies speaking of \emph{the} steady state of the forced equation. With $y_0=z_0$, it says that a change of the input by at most $\varepsilon$ changes the solution by at most $\varepsilon/\alpha$ for all time, however long. Without the lower bound on $p$ the second statement fails: for $p\equiv0$ the equation $y'=q$ adds up its inputs, and the constant input $\varepsilon$ changes the solution by $\varepsilon(t-t_0)$, which grows without bound.

\section{Exact Equations and Integrating Factors}
\label{sec:exact}

\subsection{Equations in Symmetric Form}

In this section a solution is a curve in the plane, and we use the coordinates $x$ and $y$. An equation of the form
\begin{equation}
M(x,y)+N(x,y)\,y'=0,\qquad\text{written}\qquad M\,dx+N\,dy=0 ,
\label{eq:symmetric-form}
\end{equation}
with continuous functions $M$ and $N$ on an open set $U\subseteq\mathbb R^2$, is said to be in \term{symmetric form}. A solution is a differentiable function $y(x)$ with $M(x,y(x))+N(x,y(x))\,y'(x)=0$; but the second way of writing the equation treats $x$ and $y$ alike, and a differentiable function $x(y)$ with $M(x(y),y)\,x'(y)+N(x(y),y)=0$ describes a solution curve as well. Where $N\ne0$, the equation is $y'=-M/N$, and where $M\ne0$, it is $x'=-N/M$; near points where $N=0$ the integral curves may have vertical tangents, and then they are better described as functions of $y$. Conversely, every equation $y'=f(x,y)$ can be written in symmetric form, for instance as $f\,dx-dy=0$, and in many other ways, since multiplying $M$ and $N$ by the same function without zeros does not change the solutions.

\subsection{Exact Equations}

The simplest equations in symmetric form are those whose left side is the derivative of a function along the solution.

\begin{definition}[Exact equation]
\label{def:exact}
The equation $M\,dx+N\,dy=0$ is \term{exact} on the open set $U$ if there is a function $F\in C^1(U)$ with
\[
\frac{\partial F}{\partial x}=M,\qquad\frac{\partial F}{\partial y}=N\qquad\text{on }U ,
\]
that is, if the vector field $(M,N)$ has the potential $F$ in the sense of \cref{def:potential}.
\end{definition}

\begin{theorem}[Solutions of exact equations]
\label{thm:exact-solutions}
Let $F\in C^1(U)$ with $\nabla F=(M,N)$ on the open set $U$.
\begin{enumerate}
\item A differentiable function $y\colon I\to\mathbb R$ on an interval $I$, with $(x,y(x))\in U$ for $x\in I$, is a solution of $M+Ny'=0$ if and only if $F(x,y(x))$ is constant on $I$.
\item Let $(a,b)\in U$ with $N(a,b)\ne0$. There are $r,\delta>0$ and a solution $\varphi$ on $(a-r,a+r)$ with $\varphi(a)=b$, determined implicitly by $F(x,\varphi(x))=F(a,b)$, and every solution $y$ on an interval $I\ni a$ with $y(a)=b$ and $|y(x)-b|<\delta$ for $x\in I$ coincides with $\varphi$ on $I\cap(a-r,a+r)$.
\end{enumerate}
\end{theorem}

\begin{proof}
(i) By the chain rule \eqref{eq:chain-curve} for the path $x\mapsto(x,y(x))$,
\[
\frac d{dx}F\bigl(x,y(x)\bigr)=F_x\bigl(x,y(x)\bigr)+F_y\bigl(x,y(x)\bigr)\,y'(x)=M+Ny' ,
\]
and a function on an interval is constant if and only if its derivative vanishes (\cref{cor:zero-derivative}).

(ii) Apply the implicit function theorem for one equation (\cref{thm:implicit-scalar}) to $\Phi(x,y)\coloneqq F(x,y)-F(a,b)$, which is $C^1$, vanishes at $(a,b)$ and has $\partial_y\Phi(a,b)=N(a,b)\ne0$. It gives $r,\delta>0$ and $\varphi\in C^1((a-r,a+r))$ with values in $(b-\delta,b+\delta)$ such that, for $|x-a|<r$ and $|y-b|<\delta$, $\Phi(x,y)=0$ holds exactly when $y=\varphi(x)$. In particular $\varphi(a)=b$, and $F(x,\varphi(x))=F(a,b)$ is constant, so $\varphi$ is a solution by (i). If $y$ is a solution as in the statement, then $F(x,y(x))=F(a,b)$ by (i), and for $x\in I\cap(a-r,a+r)$ the point $(x,y(x))$ lies in the box where the zeros of $\Phi$ are the points $(x,\varphi(x))$; so $y(x)=\varphi(x)$.
\end{proof}

Thus the integral curves of an exact equation are the pieces of the level curves $F(x,y)=C$ of its potential, and that through a point with $N\ne0$ passes locally only one of them; every solution with $y(a)=b$ stays within $\delta$ of $b$ near $a$, by continuity, so it coincides with $\varphi$ near $a$. Solving an exact equation therefore amounts to finding the potential; the relation $F(x,y)=C$ is the solution in \term{implicit form}, and solving it for $y$, if possible, gives the explicit form. At points where $M$ and $N$ both vanish, the critical points of $F$, the level curves may cross or degenerate, and the equation is \term{singular} there.

\subsection{The Test for Exactness}

\Cref{ch:vector-calculus} answered the question of when a vector field has a potential, and the answer translates directly.

\begin{theorem}[Test for exactness]
\label{thm:exact-test}
Let $M$ and $N$ be continuously differentiable on the open set $U\subseteq\mathbb R^2$.
\begin{enumerate}
\item If the equation $M\,dx+N\,dy=0$ is exact on $U$, then $M_y=N_x$ on $U$.
\item If $M_y=N_x$ on $U$ and $U$ is star-shaped, in particular if $U$ is a disc, an open rectangle, a half-plane or the whole plane, or more generally if $U$ is simply connected, then the equation is exact on $U$. A potential is the line integral $F(x,y)\coloneqq\int_{\boldsymbol\gamma}M\,dx+N\,dy$ along any piecewise $C^1$ path $\boldsymbol\gamma$ in $U$ from a fixed point of $U$ to $(x,y)$.
\end{enumerate}
\end{theorem}

\begin{proof}
(i) is \cref{prop:closed-condition} for the field $(M,N)$. (ii) is \cref{thm:poincare-lemma} for star-shaped sets and \cref{thm:simply-connected} for simply connected ones; that the line integral along any path is a potential is part of \cref{thm:path-independence}.
\end{proof}

On an open rectangle containing the point $(a,b)$, the line integral along the path that runs from $(a,b)$ vertically to $(a,y)$ and then horizontally to $(x,y)$ gives the explicit potential
\begin{equation}
F(x,y)=\int_b^yN(a,\eta)\,d\eta+\int_a^xM(\xi,y)\,d\xi .
\label{eq:rectangle-potential}
\end{equation}
In practice one proceeds in the same spirit: integrate $M$ with respect to $x$, which determines $F$ up to a function $c(y)$, and choose $c$ so that $F_y=N$.

The field $(2xy,\,x^2+3y^2)$ of \cref{sec:conservative} gives the exact equation
\[
2xy+(x^2+3y^2)\,y'=0 ,
\]
since $\partial_y(2xy)=2x=\partial_x(x^2+3y^2)$ on the plane. Integrating $2xy$ with respect to $x$ gives $x^2y+c(y)$, and $c'(y)=3y^2$, so $F=x^2y+y^3$, as found there with the formula of the Poincar\'e lemma. The solutions satisfy $x^2y+y^3=C$; the one through $(1,1)$ has $C=2$. Here $N=x^2+3y^2$ vanishes only at the origin, and for every fixed $x$ the function $y\mapsto x^2y+y^3$ is strictly increasing and takes every real value, so each level set $x^2y+y^3=C$ is the graph of a function $y(x)$ defined for all $x$: every solution of this equation exists on all of $\mathbb R$.

The hypothesis on the shape of $U$ matters. The equation
\[
-\frac y{x^2+y^2}\,dx+\frac x{x^2+y^2}\,dy=0
\]
on the plane without the origin comes from the angle field \eqref{eq:angle-field}, which satisfies $M_y=N_x$ but has no potential there. On the slit plane, which is star-shaped, the equation is exact with the potential $\theta=\operatorname{atan2}(y,x)$, so its integral curves are the rays $\theta=\text{constant}$ from the origin. The same conclusion can be reached directly: multiplied by $x^2+y^2$, the equation becomes $-y\,dx+x\,dy=0$, that is, $y'=y/x$, whose solutions are the lines through the origin.

\subsection{Integrating Factors}

An equation that is not exact can sometimes be made exact by multiplying it by a suitable function. A function $\mu\in C^1(U)$ without zeros is an \term{integrating factor} of $M\,dx+N\,dy=0$ if $\mu M\,dx+\mu N\,dy=0$ is exact. Since $\mu$ has no zeros, the two equations have the same solutions. On a star-shaped set the condition for $\mu$ is $(\mu M)_y=(\mu N)_x$, that is,
\[
\mu_yM-\mu_xN=\mu\,(N_x-M_y),
\]
a partial differential equation for $\mu$, which in general is no easier than the original problem. Two special cases can be solved.

\begin{proposition}[Integrating factors depending on one variable]
\label{prop:integrating-factor}
Let $M$ and $N$ be continuously differentiable on an open rectangle $U=I_1\times I_2$.
\begin{enumerate}
\item If $N$ has no zero on $U$ and $(M_y-N_x)/N=\phi(x)$ is a function of $x$ alone, then $\mu(x)\coloneqq\exp\bigl(\int\phi(x)\,dx\bigr)$, with any antiderivative of $\phi$, is an integrating factor.
\item If $M$ has no zero on $U$ and $(N_x-M_y)/M=\psi(y)$ is a function of $y$ alone, then $\mu(y)\coloneqq\exp\bigl(\int\psi(y)\,dy\bigr)$ is an integrating factor.
\end{enumerate}
Conversely, if a function $\mu(x)$ of $x$ alone, without zeros, is an integrating factor and $N$ has no zero, then $(M_y-N_x)/N=\mu'/\mu$ depends on $x$ alone.
\end{proposition}

\begin{proof}
For a $C^1$ function $\mu$ of $x$ alone, $(\mu M)_y=\mu M_y$ and $(\mu N)_x=\mu'N+\mu N_x$, so $(\mu M)_y=(\mu N)_x$ holds if and only if $\mu'N=\mu(M_y-N_x)$, that is, if and only if $\mu'/\mu=(M_y-N_x)/N$. This proves the converse, and in case (i) the function $\mu=e^\Phi$ with $\Phi'=\phi$ satisfies $\mu'/\mu=\phi$. The rectangle is star-shaped, so the equation multiplied by $\mu$ is exact by \cref{thm:exact-test}. Case (ii) is the same computation with the roles of $x$ and $y$ exchanged.
\end{proof}

The equation $(3xy+y^2)+(x^2+xy)\,y'=0$ is not exact, since $M_y=3x+2y$ and $N_x=2x+y$. But $M_y-N_x=x+y$ and $N=x(x+y)$, so $(M_y-N_x)/N=1/x$ wherever $N\ne0$, and $\mu=x$ is an integrating factor. Indeed $(3x^2y+xy^2)+(x^3+x^2y)\,y'=0$ is exact on the whole plane, with the potential $F=x^3y+x^2y^2/2$, and the solutions satisfy $x^3y+x^2y^2/2=C$. The proposition itself applies on rectangles on which $N$ has no zero, that is, which avoid the lines $x=0$ and $x+y=0$; the potential shows that the multiplied equation is in fact exact on the whole plane. The factor $x$ vanishes on the line $x=0$, where the two equations may differ.

The methods of the preceding sections are special cases. A separable equation $y'=g(x)h(y)$, written $g(x)\,dx-dy/h(y)=0$ where $h\ne0$, is exact with the potential $G(x)-H(y)$, and \cref{thm:separable}\,(iii) is \cref{thm:exact-solutions}\,(i) for this potential; the factor $1/h(y)$ is an integrating factor of $g(x)h(y)\,dx-dy=0$. A linear equation $y'+p(t)y=q(t)$, written $(p(t)y-q(t))\,dt+dy=0$, has $M_y-N_t=p(t)$ and $N=1$, so \cref{prop:integrating-factor}\,(i) gives the integrating factor $e^{\int p}$ of \cref{sec:linear-first-order}.

\subsection{Orthogonal Trajectories}

Let a family of curves be given as the level curves $F(x,y)=c$ of a $C^1$ function $F$ with $\nabla F\ne\mathbf 0$. The curves that cross every curve of the family at a right angle are its \term{orthogonal trajectories}. A curve of the family satisfies $F_x\,dx+F_y\,dy=0$, so its tangent at a point is orthogonal to $\nabla F$ there (\cref{thm:gradient-level}); the orthogonal trajectories have tangents parallel to $\nabla F$, so they satisfy
\[
F_y\,dx-F_x\,dy=0,\qquad\text{or}\qquad y'=\frac{F_y}{F_x}\quad\text{where }F_x\ne0 .
\]
They are the curves of steepest ascent of $F$, like the curve in \cref{fig:gradient-contours}. For the circles $x^2+y^2=c$, this gives $y'=y/x$, whose solutions are the lines through the origin, as expected. For the parabolas $y=cx^2$ with $x\ne0$, write $F(x,y)\coloneqq y/x^2$; then $F_x=-2y/x^3$ and $F_y=1/x^2$, so the orthogonal trajectories satisfy $y'=-x/(2y)$, a separable equation with the solutions $x^2/2+y^2=C$: the orthogonal trajectories of the parabolas are ellipses, flattened by the factor $1/\sqrt2$ (\cref{fig:orthogonal}).

\begin{figure}
\sdfloat{%
\begin{tikzpicture}
\begin{axis}[sdplot, width=0.56\textwidth, height=0.5\textwidth, xmin=-2, xmax=2, ymin=-1.6, ymax=1.6,
  axis equal image, xtick={-2,-1,0,1,2}, ytick={-1,0,1}, xlabel={$x$}, ylabel={$y$}]
\foreach \c in {0.25,0.5,1,2,-0.25,-0.5,-1,-2}{\addplot[black, line width=0.7pt, domain=-2:2, samples=100] {\c*x^2};}
\addplot[black, line width=0.7pt, domain=-2:2, samples=2] {0};
\foreach \C in {0.1,0.35,0.7,1.2,1.8}{\addplot[black!50, dashed, line width=0.7pt, domain=0:360, samples=90] ({sqrt(2*\C)*cos(x)},{sqrt(\C)*sin(x)});}
\end{axis}
\end{tikzpicture}}{The parabolas $y=cx^2$ (solid) and their orthogonal trajectories, the ellipses $x^2/2+y^2=C$ (dashed). Every ellipse crosses every parabola at a right angle.\label{fig:orthogonal}}
\end{figure}

\subsection{Conservation of Energy}

Exact equations are the mathematical form of conservation laws. A particle of unit mass that moves along a line under a force $-V'(y)$ derived from a potential energy $V$ satisfies $y''=-V'(y)$, a second-order equation. Multiplying by $y'$ gives
\[
\frac d{dt}\Bigl(\frac12y'^2+V(y)\Bigr)=y'y''+V'(y)\,y'=0 ,
\]
so the energy $E=y'^2/2+V(y)$, the sum of kinetic and potential energy, is constant along every solution. In the plane of the position $y$ and the velocity $v=y'$, the path $t\mapsto(y(t),v(t))$ therefore runs along a level curve of the function $F(y,v)\coloneqq v^2/2+V(y)$, which is the potential of the exact equation $V'(y)\,dy+v\,dv=0$; this equation is what remains of the system $y'=v$, $v'=-V'(y)$ when the time is eliminated, $dv/dy=v'/y'=-V'(y)/v$. This is the conservative force field of \cref{sec:conservative} in one dimension, with the roles of force and potential as in physics. Solving the energy equation for the velocity reduces the second-order equation to the first-order equation
\[
y'=\pm\sqrt{2\bigl(E-V(y)\bigr)},
\]
which is separable. For the pendulum, with the angle $\theta$ from the vertical, $\theta''=-(g/l)\sin\theta$, the potential is $V(\theta)=-(g/l)\cos\theta$, and the level curves of $\theta'^2/2-(g/l)\cos\theta$ are closed curves around the rest position for small energies, which describe oscillations, and open curves for large energies, which describe rotations over the top. Chapter~25 studies such pictures, the phase portraits of systems in the plane.

\section{Substitutions}
\label{sec:substitutions}

A change of the unknown function can turn an equation into one of the solvable types. This section collects the classical substitutions, a remark on scaling, and two reductions of second-order equations to first-order ones.

\subsection{Scale-Invariant Equations}

An equation of the form
\[
y'=F\Bigl(\frac yt\Bigr),
\]
with a continuous function $F$ of one variable, is traditionally called \term{homogeneous}, because its right side is homogeneous of degree $0$ in $(t,y)$; the word has nothing to do with the homogeneous linear equations of \cref{sec:linear-first-order}. The slope depends only on the ratio $y/t$, so the direction field is constant along every ray from the origin, and if $y$ is a solution, so is $\lambda y(t/\lambda)$ for every $\lambda>0$, since its derivative at $t$ is $y'(t/\lambda)=F(y(t/\lambda)/(t/\lambda))$: dilations about the origin map integral curves to integral curves. The substitution $v\coloneqq y/t$ uses this symmetry.

\begin{proposition}[Homogeneous equations]
\label{prop:homogeneous-equation}
Let $F$ be continuous on an open interval $J$, and let $I$ be an interval contained in $(0,\infty)$ or in $(-\infty,0)$. A differentiable function $y$ on $I$ with $y(t)/t\in J$ is a solution of $y'=F(y/t)$ if and only if $v(t)\coloneqq y(t)/t$ is a solution of the separable equation
\[
v'=\frac{F(v)-v}t .
\]
\end{proposition}

\begin{proof}
From $y=tv$ we get $y'=v+tv'$, so $y'=F(v)$ if and only if $tv'=F(v)-v$.
\end{proof}

The zeros of $F(v)-v$ give the solutions $y=vt$ along rays. For the equation
\[
y'=\frac{t^2+y^2}{2ty}\qquad(t>0,\ y>0),
\]
$F(v)=(1+v^2)/(2v)$ and $F(v)-v=(1-v^2)/(2v)$, which vanishes at $v=1$; so $y=t$ is a solution. For $0<v<1$, separation of variables gives $\int2v\,dv/(1-v^2)=\int dt/t$, that is, $-\ln(1-v^2)=\ln t+c$, or $1-v^2=C/t$ with $C>0$, and for $v>1$ the same computation gives this relation with $C<0$. So the solutions are $y=\sqrt{t^2-Ct}$, on the interval of those $t>0$ where $t^2-Ct>0$, together with $y=t$ for $C=0$; their graphs are arcs of the hyperbolas $(t-C/2)^2-y^2=C^2/4$, a family invariant under dilations, as it must be.

\subsection{Bernoulli Equations}

\begin{proposition}[Bernoulli equations]
\label{prop:bernoulli}
Let $p$ and $q$ be continuous on an interval $I$, and let $n$ be a real number other than $0$ and $1$. A differentiable function $y\colon I\to(0,\infty)$ is a solution of
\[
y'+p(t)\,y=q(t)\,y^n
\]
if and only if $v\coloneqq y^{1-n}$ is a solution of the linear equation $v'+(1-n)\,p(t)\,v=(1-n)\,q(t)$.
\end{proposition}

\begin{proof}
By the chain rule $v'=(1-n)y^{-n}y'$, and $y^{-n}(y'+py-qy^n)=y'y^{-n}+pv-q=\bigl(v'+(1-n)pv-(1-n)q\bigr)/(1-n)$. Since $y^{-n}>0$, the left side vanishes exactly when the right side does.
\end{proof}

The logistic equation $y'-ry=-(r/K)y^2$ is a Bernoulli equation with $n=2$. The substitution $v=1/y$ gives $v'+rv=r/K$, whose solutions are $v=1/K+(v_0-1/K)e^{-r(t-t_0)}$ by \eqref{eq:duhamel}; taking reciprocals recovers \eqref{eq:logistic-solution}, now without partial fractions. The reciprocal of a logistic population relaxes exponentially to its equilibrium value $1/K$, as the excess temperature does under Newton's law of cooling.

\subsection{Riccati Equations}

An equation whose right side is a quadratic polynomial in $y$,
\begin{equation}
y'=a(t)+b(t)\,y+c(t)\,y^2 ,
\label{eq:riccati}
\end{equation}
with continuous coefficients, is a \term{Riccati equation}. For $c\equiv0$ it is linear, and for $a\equiv0$ it is a Bernoulli equation; in general it cannot be solved by quadratures, as the note in \cref{sec:direction-fields} explained. But one solution suffices to find all the others.

\begin{proposition}[Riccati equations]
\label{prop:riccati}
Let $y_1$ be a solution of \eqref{eq:riccati} on an interval $I$. A differentiable function $y$ on $I$ is a solution if and only if $u\coloneqq y-y_1$ satisfies
\[
u'=\bigl(b(t)+2c(t)y_1(t)\bigr)\,u+c(t)\,u^2 ,
\]
and on an interval where $u$ has no zero this holds if and only if $v\coloneqq1/u$ satisfies the linear equation $v'=-\bigl(b(t)+2c(t)y_1(t)\bigr)\,v-c(t)$.
\end{proposition}

\begin{proof}
Expanding, $a+b(y_1+u)+c(y_1+u)^2=(a+by_1+cy_1^2)+(b+2cy_1)u+cu^2$, and the first bracket is $y_1'$. So $y'=y_1'+u'$ equals the right side of \eqref{eq:riccati} at $y$ exactly when $u'=(b+2cy_1)u+cu^2$. If $u$ has no zero, $v'=-u'/u^2$, and dividing the equation for $u$ by $-u^2$ gives the equation for $v$, and conversely.
\end{proof}

The equation $y'=y^2-2/t^2$ on $t>0$ has the solution $y_1=1/t$, as one checks: $-1/t^2=1/t^2-2/t^2$. Here $b\equiv0$ and $c\equiv1$, so $v$ satisfies $v'=-(2/t)v-1$. The integrating factor $t^2$ gives $(t^2v)'=-t^2$, so $t^2v=(C-t^3)/3$, and
\[
y=\frac1t+\frac{3t^2}{C-t^3} ,
\]
which for $C=0$ is the second simple solution $-2/t$. A Riccati equation is also related to a linear equation of the second order: if $c\equiv1$ and $b\equiv0$, the substitution $y=-w'/w$ gives $y'=-w''/w+(w'/w)^2$, and $y'=a+y^2$ becomes $w''+a(t)\,w=0$. For $y'=y^2-2/t^2$, the linear equation $w''=2w/t^2$ has the solutions $t^2$ and $1/t$, which give back $y=-2/t$ and $y=1/t$. Conversely, a first-order nonlinear equation is sometimes best solved through a linear equation of higher order, which \cref{ch:linear-higher-order} and Chapter~26 treat.

\subsection{Scaling}

The constants in a model can often be absorbed into the units of measurement. In the logistic equation \eqref{eq:logistic}, measure the population in units of the carrying capacity and the time in units of $1/r$: with $u\coloneqq y/K$ and $s\coloneqq rt$, the chain rule gives
\[
\frac{du}{ds}=\frac1{rK}\,y'=u(1-u) .
\]
A single equation without constants describes all logistic populations, and the solutions of \cref{fig:logistic} are those of every logistic equation, once the axes are labelled in these units. In the same way, the falling body with quadratic resistance becomes $du/ds=1-u^2$ with $u\coloneqq v/v_{\mathrm T}$ and $s\coloneqq gt/v_{\mathrm T}$, whose solution from rest is $u=\tanh s$. The new variables are dimensionless: they are ratios of quantities measured in the same units, and the equations relate them without reference to metres, seconds or individuals. When a model has more constants than can be absorbed, scaling reduces them to a few dimensionless combinations, and the behaviour of the solutions depends only on these. The logistic equation with harvesting of \cref{sec:autonomous}, $y'=ry(1-y/K)-h$, becomes $u'=u(1-u)-\eta$ with the single parameter $\eta\coloneqq h/(rK)$, the harvest as a fraction of $rK$.

\subsection{Second-Order Equations of the First Order in Disguise}

Two kinds of second-order equations reduce to first-order equations. If the unknown function itself is absent, $y''=f(t,y')$, the derivative $p\coloneqq y'$ satisfies the first-order equation $p'=f(t,p)$, and $y$ is recovered by one integration. The catenary equation \eqref{eq:catenary} was solved in this way in \cref{subsec:catenary}, and so is the motion of a falling body, whose position $y$ satisfies $y''=g-cy'$ or $y''=g-cy'^2$.

If the independent variable is absent, $y''=f(y,y')$, the velocity can be regarded as a function of the position rather than of the time.

\begin{proposition}[Equations without the independent variable]
\label{prop:reduction-autonomous}
Let $y$ be a twice differentiable solution of $y''=f(y,y')$ on an interval $I$ with $y'(t)\ne0$ for all $t\in I$. Then $y$ is strictly monotone, and the function $p$ on the interval $y(I)$ defined by $p(y(t))\coloneqq y'(t)$ satisfies the first-order equation
\[
p\,\frac{dp}{dy}=f(y,p) .
\]
\end{proposition}

\begin{proof}
The derivative $y'$ is continuous, being differentiable, and has no zero, so it has a constant sign and $y$ is strictly monotone (\cref{cor:ivt,cor:monotone-derivative}). Let $\tau\colon y(I)\to I$ be the inverse of $y$; by \cref{thm:inverse-derivative} it is differentiable with $\tau'(\eta)=1/y'(\tau(\eta))$. Then $p=y'\circ\tau$, and by the chain rule
\[
p'(\eta)=y''\bigl(\tau(\eta)\bigr)\,\tau'(\eta)=\frac{f\bigl(\eta,p(\eta)\bigr)}{p(\eta)} .\qedhere
\]
\end{proof}

Once $p$ is known, $y$ is found from the autonomous first-order equation $y'=p(y)$. When $f$ does not depend on $y'$, the equation $p\,dp/dy=f(y)$ integrates to $p^2/2-\int f=\text{constant}$, the conservation of energy of \cref{sec:exact}. A body launched vertically from the surface of a planet of mass $M$ and radius $R$ with the speed $v_0$ has the distance $r(t)$ from the centre with $r''=-GM/r^2$. While it rises, $p(r)\coloneqq r'$ satisfies $p\,dp/dr=-GM/r^2$, so that
\[
\frac{p^2}2-\frac{GM}r=\frac{v_0^2}2-\frac{GM}R .
\]
If $v_0^2<2GM/R$, the speed vanishes at the height where $1/r=1/R-v_0^2/(2GM)$, and the body falls back. If $v_0^2\ge2GM/R$, then $p^2\ge2GM/r>0$ for all $r\ge R$, so the body never stops rising, and since $r'\ge\sqrt{2GM/r}$, the function $r^{3/2}$ grows at least linearly in $t$ and $r(t)\to\infty$: the body escapes. The threshold $\sqrt{2GM/R}$ is the \term{escape velocity}; its kinetic energy per unit mass, $GM/R$, is the energy that \cref{sec:conservative} found from the work done by gravity to be needed for escape.

\section{Autonomous Equations and Stability}
\label{sec:autonomous}

\subsection{Equilibria and the Phase Line}

An autonomous equation
\begin{equation}
y'=f(y),
\label{eq:autonomous}
\end{equation}
with $f$ continuous on an open interval $J$, describes a quantity whose rate of change depends only on its present value, not on the time at which it has this value. Two consequences are immediate. If $y(t)$ is a solution, so is $y(t-c)$ for every constant $c$, since its derivative at $t$ is $y'(t-c)=f(y(t-c))$: shifting an integral curve horizontally gives another integral curve, as \cref{fig:direction-fields}B showed. And every zero $y^*$ of $f$ gives a constant solution $y\equiv y^*$, an \term{equilibrium} of the equation, also called a stationary or critical point. Between the equilibria, the sign of $f$ tells in which direction the solutions move: upwards where $f>0$ and downwards where $f<0$. Marking the equilibria on a line and drawing arrows between them gives the \term{phase line} of the equation, a complete qualitative picture of its solutions, as the following results show (\cref{fig:allee}A).

\begin{theorem}[Solutions between equilibria]
\label{thm:between-equilibria}
Let $f$ be continuous on an open interval $J$, and let $(a,b)\subseteq J$, with $-\infty\le a<b\le\infty$, be an interval on which $f>0$, whose ends are zeros of $f$ or ends of $J$. Let $t_0\in\mathbb R$ and $y_0\in(a,b)$, and let
\[
H(y)\coloneqq\int_{y_0}^y\frac{d\eta}{f(\eta)}\quad(a<y<b),
\]
and let $T_+\coloneqq\lim_{y\to b^-}H(y)\in(0,\infty]$ and $T_-\coloneqq\lim_{y\to a^+}H(y)\in[-\infty,0)$.
\begin{enumerate}
\item The function $y(t)\coloneqq H^{-1}(t-t_0)$ is a strictly increasing solution of $y'=f(y)$, $y(t_0)=y_0$, on the interval $(t_0+T_-,t_0+T_+)$, and $y(t)\to b$ as $t\to t_0+T_+$ and $y(t)\to a$ as $t\to t_0+T_-$.
\item Every solution of this initial value problem with values in $(a,b)$ is a restriction of the solution in (i).
\item If $b\in J$ and there are $L>0$ and $\varepsilon>0$ with $f(\eta)\le L(b-\eta)$ for $b-\varepsilon<\eta<b$, in particular if $f$ is differentiable at the zero $b$, then $T_+=\infty$: the solution approaches the equilibrium $b$ but does not reach it in finite time. The same holds at $a$ if $f(\eta)\le L(\eta-a)$ near $a$.
\item If $b=\infty$, then $T_+=\int_{y_0}^\infty d\eta/f(\eta)$, and the solution tends to $\infty$ in finite time if and only if this improper integral converges.
\end{enumerate}
The same holds on an interval on which $f<0$, with \emph{decreasing} in place of \emph{increasing} and the roles of $a$ and $b$ exchanged.
\end{theorem}

\begin{proof}
(i) and (ii) are \cref{thm:separable} with $g\equiv1$, $G(t)=t-t_0$ and $h=f$, for which $K_0=(a,b)$. Since $H'=1/f>0$, $H$ is strictly increasing, so the limits $T_\pm$ exist in the extended sense, as for monotone sequences (\cref{thm:monotone-one-sided}), and $H$ maps $(a,b)$ onto $(T_-,T_+)$. The set of $t$ with $G(t)\in H(K_0)$ is therefore the interval $(t_0+T_-,t_0+T_+)$. The inverse $H^{-1}$ is a strictly increasing bijection of $(T_-,T_+)$ onto $(a,b)$; so for every $c<b$ it exceeds $c$ on $(H(c),T_+)$, which shows that $y(t)\to b$ as $t\to t_0+T_+$, and similarly at the other end.

(iii) Let $y_1\coloneqq\max(y_0,b-\varepsilon/2)$. For $y_1\le y<b$,
\[
H(y)=H(y_1)+\int_{y_1}^y\frac{d\eta}{f(\eta)}\ge H(y_1)+\frac1L\int_{y_1}^y\frac{d\eta}{b-\eta}=H(y_1)+\frac1L\ln\frac{b-y_1}{b-y},
\]
which tends to $\infty$ as $y\to b^-$. If $f$ is differentiable at $b$ with $f(b)=0$, then $f(\eta)=f'(b)(\eta-b)+o(\eta-b)$, so $f(\eta)\le(|f'(b)|+1)(b-\eta)$ for $\eta<b$ close to $b$. The argument at $a$ is the same.

(iv) is the definition of the improper integral (\cref{def:improper}).
\end{proof}

When solutions are unique, these are all the solutions.

\begin{proposition}[Monotonicity of solutions]
\label{prop:autonomous-monotone}
Let $f$ be continuous on $J$, and suppose that every initial value problem for $y'=f(y)$ has at most one solution on each interval, as is the case if $f$ is continuously differentiable (\cref{cor:local-uniqueness}). Then every solution on an interval is either an equilibrium or strictly monotone, and in the second case its values lie in an interval between two consecutive zeros of $f$, or between a zero and an end of $J$, so that \cref{thm:between-equilibria} describes it.
\end{proposition}

\begin{proof}
If $f(y(t_1))=0$ for some $t_1$, the constant function $y(t_1)$ is a solution with the same value at $t_1$, and by uniqueness $y$ is this constant. Otherwise $f(y(t))$ never vanishes; it is continuous in $t$, so it has a constant sign (\cref{cor:ivt}), and so does $y'$, and $y$ is strictly monotone. Its values form an interval (\cref{cor:interval-image}) that contains no zero of $f$, which is therefore contained in one of the intervals of the theorem.
\end{proof}

Without uniqueness, the proposition fails: a solution of $y'=\sqrt{|y|}$ may rest at the equilibrium $0$ and leave it later, as in \cref{fig:two-warnings}B. In that example the condition of (iii) fails at the zero, and $\int_0^1d\eta/\sqrt\eta$ converges, so that the solutions reach the equilibrium in finite time, backwards; the leaking tank of \cref{sec:separable} reaches its equilibrium forwards. For a continuously differentiable $f$ with isolated zeros, the picture is simple. Each nonconstant solution moves monotonically, in the direction of the arrows of the phase line, between two neighbouring points each of which is an equilibrium or an end of $J$. It approaches a neighbouring equilibrium without reaching it in finite time, by (iii); when $J=\mathbb R$ and no equilibrium lies ahead, it escapes to $\pm\infty$, in finite time exactly when the integral in (iv) converges.

\subsection{Stability}

An equilibrium matters in practice only if it is robust: a population, a temperature or a chemical concentration is never exactly at its equilibrium value, and the question is whether small deviations die out or grow.

\begin{definition}[Stability]
\label{def:stability}
Let $f$ be continuously differentiable on $J$ and $f(y^*)=0$. The equilibrium $y^*$ of $y'=f(y)$ is \term{stable} if for every $\varepsilon>0$ there is $\delta>0$ such that for $|y_0-y^*|<\delta$ the solution with $y(t_0)=y_0$ exists for all $t\ge t_0$ and satisfies $|y(t)-y^*|<\varepsilon$ for all $t\ge t_0$. It is \term{asymptotically stable} if moreover $\delta$ can be chosen so that $y(t)\to y^*$ as $t\to\infty$ for $|y_0-y^*|<\delta$. It is \term{unstable} if it is not stable.
\end{definition}

The derivative of $f$ at the equilibrium decides the question in most cases. Near $y^*$, $f(y)\approx f'(y^*)(y-y^*)$, so the deviation $u=y-y^*$ approximately satisfies the \term{linearised equation} $u'=f'(y^*)u$, whose solutions decay if $f'(y^*)<0$ and grow if $f'(y^*)>0$.

\begin{theorem}[Linearised stability]
\label{thm:linearised-stability}
Let $f$ be continuously differentiable on $J$ and $f(y^*)=0$.
\begin{enumerate}
\item If $f'(y^*)<0$, then $y^*$ is asymptotically stable. More precisely, for every $\alpha$ with $0<\alpha<-f'(y^*)$ there is $\delta>0$ such that, for $|y_0-y^*|\le\delta$, the solution with $y(t_0)=y_0$ satisfies
\[
|y(t)-y^*|\le|y_0-y^*|\,e^{-\alpha(t-t_0)}\qquad(t\ge t_0).
\]
\item If $f'(y^*)>0$, then $y^*$ is unstable.
\end{enumerate}
\end{theorem}

\begin{proof}
(i) Since $f'$ is continuous, there is $\delta>0$ with $[y^*-\delta,y^*+\delta]\subseteq J$ and $f'(\eta)\le-\alpha$ for $|\eta-y^*|\le\delta$. By the mean value theorem, $f(\eta)=f'(\xi)(\eta-y^*)$ with some $\xi$ between $y^*$ and $\eta$, so $f<0$ on $(y^*,y^*+\delta]$ and $f>0$ on $[y^*-\delta,y^*)$. Let $y^*<y_0\le y^*+\delta$; the case $y_0<y^*$ is symmetric, and for $y_0=y^*$ there is nothing to prove. By \cref{thm:between-equilibria}, applied to the interval of negative values of $f$ above $y^*$, the solution decreases towards $y^*$, and since $f$ is differentiable at $y^*$ it exists for all $t\ge t_0$; so its values lie in $(y^*,y_0]$. The function $w\coloneqq(y-y^*)^2$ satisfies
\[
w'=2(y-y^*)f(y)=2f'(\xi_t)(y-y^*)^2\le-2\alpha w ,
\]
with $\xi_t$ between $y^*$ and $y(t)$, so $(e^{2\alpha t}w)'=e^{2\alpha t}(w'+2\alpha w)\le0$, and $e^{2\alpha t}w(t)$ does not increase (\cref{cor:monotone-derivative}). Hence $w(t)\le w(t_0)e^{-2\alpha(t-t_0)}$, which is the claim.

(ii) Choose $\alpha>0$ and $\delta>0$ with $[y^*,y^*+\delta]\subseteq J$ and $f'(\eta)\ge\alpha$ for $y^*\le\eta\le y^*+\delta$; then $f(\eta)\ge\alpha(\eta-y^*)>0$ for $y^*<\eta\le y^*+\delta$. Let $y^*<y_0\le y^*+\delta$, and suppose that the solution with $y(t_0)=y_0$ exists and stays in $(y^*-\delta,y^*+\delta)$ for all $t\ge t_0$. It increases, so $y^*<y(t)<y^*+\delta$, and $u\coloneqq y-y^*$ satisfies $u'=f(y)\ge\alpha u$. As in (i), $e^{-\alpha t}u(t)$ does not decrease, so $u(t)\ge u(t_0)e^{\alpha(t-t_0)}$, which exceeds $\delta$ for large $t$, a contradiction. So for $\varepsilon=\delta$, initial values arbitrarily close to $y^*$ lead to solutions that leave the $\varepsilon$-neighbourhood, and $y^*$ is not stable.
\end{proof}

The theorem says nothing when $f'(y^*)=0$, and then anything can happen. The equation $y'=-y^3$ has the asymptotically stable equilibrium $0$, but its solutions $y=y_0/\sqrt{1+2y_0^2t}$ decay only like $1/\sqrt{2t}$, not exponentially; $y'=y^3$ has the unstable equilibrium $0$, and its solutions blow up; and $y'=y^2$ has an equilibrium that attracts solutions from below and repels them above, so it is unstable. The number $-f'(y^*)$ is the rate at which deviations decay, and its reciprocal $-1/f'(y^*)$ is the time scale on which the equilibrium recovers from a disturbance; a small value of $|f'(y^*)|$ means a slow recovery even when the equilibrium is stable.

For the logistic equation, $f(y)=ry(1-y/K)$ has $f'(y)=r(1-2y/K)$, so $f'(0)=r>0$ and $f'(K)=-r<0$: the extinct population is an unstable equilibrium, since a few individuals grow, and the carrying capacity is asymptotically stable, with the recovery rate $r$.

\subsection{Thresholds}

Some species cannot persist at low densities, for instance because individuals fail to find mates. A simple model with this property is
\[
y'=ry\Bigl(\frac yA-1\Bigr)\Bigl(1-\frac yK\Bigr),\qquad0<A<K,
\]
which has the equilibria $0$, $A$ and $K$~\cite[Sect.~2.5]{boyce2017}. The derivative of the right side is $-r$ at $0$, $r(1-A/K)>0$ at $A$ and $-r(K/A-1)<0$ at $K$. By \cref{thm:linearised-stability}, extinction and the carrying capacity are both asymptotically stable, and the threshold $A$ between them is unstable. A population above $A$ grows to $K$; one below $A$ dies out. \Cref{fig:allee} shows the phase line and the solutions for $r=K=1$ and $A=0.3$. An equation with two stable equilibria is called \term{bistable}; which equilibrium a solution approaches depends on which side of the threshold it starts, and a disturbance that pushes the population across the threshold changes its fate permanently.

\begin{figure}
\sdpanel{A}{\begin{tikzpicture}
\begin{axis}[sdplot, width=0.47\textwidth, height=0.42\textwidth, xmin=-0.08, xmax=1.3, ymin=-0.62, ymax=0.62,
  xtick={0,0.3,1}, xticklabels={0,$A$,$K$}, ytick={-0.5,0,0.5}, xlabel={$y$}, ylabel={$f(y)$}]
\addplot[black!45, line width=0.4pt, domain=-0.08:1.3, samples=2] {0};
\addplot[black, line width=0.8pt, domain=-0.08:1.3, samples=120] {x*(x/0.3-1)*(1-x)};
\draw[black, line width=0.7pt, -{Latex[length=1.8mm]}] (axis cs:0.24,0) -- (axis cs:0.06,0);
\draw[black, line width=0.7pt, -{Latex[length=1.8mm]}] (axis cs:0.36,0) -- (axis cs:0.94,0);
\draw[black, line width=0.7pt, -{Latex[length=1.8mm]}] (axis cs:1.26,0) -- (axis cs:1.06,0);
\fill[black] (axis cs:0,0) circle (1.8pt) (axis cs:1,0) circle (1.8pt);
\draw[black, line width=0.6pt, fill=white] (axis cs:0.3,0) circle (1.8pt);
\end{axis}
\end{tikzpicture}}\hfill
\sdpanel{B}{\begin{tikzpicture}
\begin{axis}[sdplot, width=0.47\textwidth, height=0.42\textwidth, xmin=0, xmax=8, ymin=0, ymax=1.35,
  xtick={0,2,4,6,8}, ytick={0,0.3,1}, yticklabels={0,$A$,$K$}, xlabel={$t$}, ylabel={$y$}]
\addplot[black!45, dashed, line width=0.6pt, domain=0:8, samples=2] {0.3};
\addplot[black!45, dashed, line width=0.6pt, domain=0:8, samples=2] {1};
\addplot[black, line width=0.7pt, smooth] coordinates {(0.00,0.050) (0.25,0.041) (0.50,0.033) (0.75,0.027) (1.00,0.021) (1.25,0.017) (1.50,0.013) (1.75,0.011) (2.00,0.008) (2.25,0.007) (2.50,0.005) (2.75,0.004) (3.00,0.003) (3.25,0.002) (3.50,0.002) (3.75,0.001) (4.00,0.001) (4.25,0.001) (4.50,0.001) (4.75,0.001) (5.00,0.000) (5.25,0.000) (5.50,0.000) (5.75,0.000) (6.00,0.000) (6.25,0.000) (6.50,0.000) (6.75,0.000) (7.00,0.000) (7.25,0.000) (7.50,0.000) (7.75,0.000)};
\addplot[black, line width=0.7pt, smooth] coordinates {(0.00,0.200) (0.25,0.186) (0.50,0.171) (0.75,0.156) (1.00,0.140) (1.25,0.124) (1.50,0.108) (1.75,0.093) (2.00,0.079) (2.25,0.066) (2.50,0.055) (2.75,0.045) (3.00,0.037) (3.25,0.030) (3.50,0.024) (3.75,0.019) (4.00,0.015) (4.25,0.012) (4.50,0.009) (4.75,0.007) (5.00,0.006) (5.25,0.004) (5.50,0.004) (5.75,0.003) (6.00,0.002) (6.25,0.002) (6.50,0.001) (6.75,0.001) (7.00,0.001) (7.25,0.001) (7.50,0.000) (7.75,0.000)};
\addplot[black, line width=0.7pt, smooth] coordinates {(0.00,0.270) (0.25,0.265) (0.50,0.259) (0.75,0.251) (1.00,0.243) (1.25,0.234) (1.50,0.223) (1.75,0.212) (2.00,0.199) (2.25,0.185) (2.50,0.170) (2.75,0.155) (3.00,0.139) (3.25,0.123) (3.50,0.107) (3.75,0.092) (4.00,0.078) (4.25,0.065) (4.50,0.054) (4.75,0.044) (5.00,0.036) (5.25,0.029) (5.50,0.023) (5.75,0.019) (6.00,0.015) (6.25,0.012) (6.50,0.009) (6.75,0.007) (7.00,0.006) (7.25,0.004) (7.50,0.003) (7.75,0.003)};
\addplot[black, line width=0.7pt, smooth] coordinates {(0.00,0.330) (0.25,0.336) (0.50,0.344) (0.75,0.353) (1.00,0.364) (1.25,0.377) (1.50,0.394) (1.75,0.415) (2.00,0.441) (2.25,0.474) (2.50,0.514) (2.75,0.563) (3.00,0.622) (3.25,0.688) (3.50,0.758) (3.75,0.826) (4.00,0.883) (4.25,0.926) (4.50,0.955) (4.75,0.974) (5.00,0.985) (5.25,0.991) (5.50,0.995) (5.75,0.997) (6.00,0.998) (6.25,0.999) (6.50,1.000) (6.75,1.000) (7.00,1.000) (7.25,1.000) (7.50,1.000) (7.75,1.000)};
\addplot[black, line width=0.7pt, smooth] coordinates {(0.00,0.450) (0.25,0.485) (0.50,0.527) (0.75,0.579) (1.00,0.640) (1.25,0.709) (1.50,0.779) (1.75,0.844) (2.00,0.897) (2.25,0.936) (2.50,0.961) (2.75,0.978) (3.00,0.987) (3.25,0.993) (3.50,0.996) (3.75,0.998) (4.00,0.999) (4.25,0.999) (4.50,1.000) (4.75,1.000) (5.00,1.000) (5.25,1.000) (5.50,1.000) (5.75,1.000) (6.00,1.000) (6.25,1.000) (6.50,1.000) (6.75,1.000) (7.00,1.000) (7.25,1.000) (7.50,1.000) (7.75,1.000)};
\addplot[black, line width=0.7pt, smooth] coordinates {(0.00,0.750) (0.25,0.818) (0.50,0.876) (0.75,0.921) (1.00,0.952) (1.25,0.972) (1.50,0.984) (1.75,0.991) (2.00,0.995) (2.25,0.997) (2.50,0.998) (2.75,0.999) (3.00,0.999) (3.25,1.000) (3.50,1.000) (3.75,1.000) (4.00,1.000) (4.25,1.000) (4.50,1.000) (4.75,1.000) (5.00,1.000) (5.25,1.000) (5.50,1.000) (5.75,1.000) (6.00,1.000) (6.25,1.000) (6.50,1.000) (6.75,1.000) (7.00,1.000) (7.25,1.000) (7.50,1.000) (7.75,1.000)};
\addplot[black, line width=0.7pt, smooth] coordinates {(0.00,1.300) (0.25,1.123) (0.50,1.060) (0.75,1.032) (1.00,1.017) (1.25,1.009) (1.50,1.005) (1.75,1.003) (2.00,1.002) (2.25,1.001) (2.50,1.000) (2.75,1.000) (3.00,1.000) (3.25,1.000) (3.50,1.000) (3.75,1.000) (4.00,1.000) (4.25,1.000) (4.50,1.000) (4.75,1.000) (5.00,1.000) (5.25,1.000) (5.50,1.000) (5.75,1.000) (6.00,1.000) (6.25,1.000) (6.50,1.000) (6.75,1.000) (7.00,1.000) (7.25,1.000) (7.50,1.000) (7.75,1.000)};
\end{axis}
\end{tikzpicture}}
\caption{A population with a threshold, $y'=y(y/A-1)(1-y)$ with $A=0.3$. (A) The graph of the right side $f$ and the phase line on the horizontal axis: the arrows show the direction of motion, the stable equilibria $0$ and $K=1$ are filled, the unstable threshold $A$ is open. (B) Solutions for several initial values; the dashed lines mark the equilibria $A$ and $K$. Solutions that start above the threshold approach $K$, and those that start below it die out.\label{fig:allee}}
\end{figure}

\subsection{Harvesting and a Bifurcation}

Let a logistic population be harvested at a constant rate $h\ge0$, for instance by fishing:
\[
y'=ry\Bigl(1-\frac yK\Bigr)-h .
\]
By the scaling of \cref{sec:substitutions}, it suffices to study $u'=u(1-u)-\eta$ with $\eta=h/(rK)$. The equilibria are the roots of $u^2-u+\eta=0$,
\[
u_\pm=\frac{1\pm\sqrt{1-4\eta}}2\qquad(0\le\eta\le1/4),
\]
and the derivative of the right side, $1-2u$, is $\mp\sqrt{1-4\eta}$ at $u_\pm$. For $0<\eta<1/4$, the upper equilibrium $u_+$ is asymptotically stable and the lower one $u_-$ unstable: a harvested population settles at $u_+$, provided it never falls below $u_-$, below which it collapses. As the harvest grows, the two equilibria approach each other, the recovery rate $\sqrt{1-4\eta}$ of the stable one tends to $0$, and at $\eta=1/4$ they merge at $u=1/2$. For $\eta>1/4$ there is no equilibrium at all, since $u(1-u)-\eta\le1/4-\eta<0$ for all $u$: every solution decreases, and by \cref{thm:between-equilibria} it reaches $0$ in finite time, because the right side is bounded away from zero. Writing $\eta-u(1-u)=(u-1/2)^2+\beta^2$ with $\beta\coloneqq\sqrt{\eta-1/4}$, the time it takes to fall from $u=1$ to $u=0$ is
\[
\int_0^1\frac{du}{(u-1/2)^2+\beta^2}=\frac2\beta\arctan\frac1{2\beta} .
\]
For $\eta=0.3$ this is $10.3$, in units of $1/r$; as $\eta$ decreases to $1/4$, it grows like $\pi/\beta$, because the solution lingers near $u=1/2$, where the right side is smallest.

\Cref{fig:bifurcation} shows the equilibria as functions of the parameter. A qualitative change of the solutions at a critical value of a parameter is a \term{bifurcation}, and the merging and disappearance of a stable and an unstable equilibrium is called a \term{saddle-node bifurcation}~\cite[Sect.~3.1]{strogatz2015}. The critical harvest $h=rK/4$ is the largest rate that the population can sustain, the maximum of the natural growth $ry(1-y/K)$, attained at $y=K/2$. The model shows what makes harvesting at this rate hazardous: just below it, the stable equilibrium recovers very slowly from disturbances, and just above it, the population collapses, however large it is.

\begin{figure}
\sdfloat{%
\begin{tikzpicture}
\begin{axis}[sdplot, width=0.56\textwidth, height=0.42\textwidth, xmin=0, xmax=0.4, ymin=0, ymax=1.08,
  xtick={0,0.1,0.2,0.25,0.4}, xticklabels={0,0.1,0.2,$1/4$,0.4}, ytick={0,0.5,1}, yticklabels={0,$1/2$,1}, xlabel={$\eta$}, ylabel={$u$}]
\addplot[black, line width=0.9pt, domain=0.5:1, samples=60] ({x*(1-x)},{x});
\addplot[black, dashed, line width=0.9pt, domain=0:0.5, samples=60] ({x*(1-x)},{x});
\pgfplotsinvokeforeach{0.1,0.2}{
  \draw[black!50, line width=0.5pt, -{Latex[length=1.5mm]}] (axis cs:#1,1.05) -- (axis cs:#1,{(1+sqrt(1-4*#1))/2+0.04});
  \draw[black!50, line width=0.5pt, -{Latex[length=1.5mm]}] (axis cs:#1,{(1-sqrt(1-4*#1))/2+0.04}) -- (axis cs:#1,{(1+sqrt(1-4*#1))/2-0.04});
  \draw[black!50, line width=0.5pt, -{Latex[length=1.5mm]}] (axis cs:#1,{(1-sqrt(1-4*#1))/2-0.03}) -- (axis cs:#1,0.01);}
\draw[black!50, line width=0.5pt, -{Latex[length=1.5mm]}] (axis cs:0.33,1.05) -- (axis cs:0.33,0.01);
\fill[black] (axis cs:0.25,0.5) circle (1.4pt);
\end{axis}
\end{tikzpicture}}{The saddle-node bifurcation of the harvested logistic equation $u'=u(1-u)-\eta$. For $\eta<1/4$ there are a stable equilibrium (solid) and an unstable one (dashed); the grey arrows show the direction of motion of the solutions. They merge at $\eta=1/4$ (dot), and for $\eta>1/4$ every solution decreases and reaches $0$ in finite time.\label{fig:bifurcation}}
\end{figure}

\subsection{Blow-up and Potentials}

Part (iv) of \cref{thm:between-equilibria} decides whether solutions escape to infinity in finite time. For $y'=1+y^2$, the integral $\int_{y_0}^\infty d\eta/(1+\eta^2)=\pi/2-\arctan y_0$ converges, and the solution $y=\tan(t-t_0+\arctan y_0)$ blows up at the time $t_0+\pi/2-\arctan y_0$. For $y'=y^p$ with $p>1$ and $y_0>0$ the integral converges as well, by \cref{prop:p-integrals}, while for $p\le1$ it diverges and the solutions exist for all $t\ge t_0$. Growth faster than exponential does not by itself imply blow-up: for $y'=y\ln y$ with $y_0>1$, the integral $\int_{y_0}^\infty d\eta/(\eta\ln\eta)$ diverges, since $\ln\ln\eta$ is an antiderivative, and the solution $y=\exp(e^{t-t_0}\ln y_0)$ exists for all time, growing doubly exponentially.

Every continuous $f$ on $J$ has an antiderivative, so every autonomous equation can be written as $y'=-V'(y)$ with a \term{potential} $V\coloneqq-\int f$. Along a solution,
\[
\frac d{dt}V\bigl(y(t)\bigr)=V'(y)\,y'=-f(y)^2\le0 ,
\]
with equality only at equilibria, so every solution runs downhill on the graph of $V$, like a ball in a very viscous liquid that has no inertia. The equilibria are the critical points of $V$, and since $f'=-V''$, \cref{thm:linearised-stability} says that strict local minima with $V''(y^*)>0$ are asymptotically stable and strict local maxima with $V''(y^*)<0$ are unstable. For the logistic equation, $V(y)=-ry^2/2+ry^3/(3K)$ has its local maximum at $0$ and its local minimum at $K$. Chapter~25 extends the phase line to the phase plane of an autonomous system of two equations, where solutions can also circle around equilibria, and where the linearisation at an equilibrium is a $2\times2$ matrix, whose eigenvalues take the place of $f'(y^*)$.

\section{Existence and Uniqueness}
\label{sec:existence-uniqueness}

\subsection{Existence}

The examples of \cref{sec:ode-basics} raised two questions: when does an initial value problem have a solution, and when only one? The first has a satisfying answer: continuity of the right side suffices, at least for a short time.

\begin{theorem}[Peano's existence theorem]
\label{thm:peano}
Let $f$ be continuous on the rectangle $R\coloneqq[t_0-a,t_0+a]\times[y_0-b,y_0+b]$ with $a,b>0$, let $M\coloneqq\max_R|f|$, and let $\alpha\coloneqq\min(a,b/M)$, or $\alpha\coloneqq a$ if $M=0$. Then the initial value problem $y'=f(t,y)$, $y(t_0)=y_0$, has at least one solution on $[t_0-\alpha,t_0+\alpha]$.
\end{theorem}

\begin{proof}[sketch]
For a step $h>0$, the polygon of Euler's method (\cref{sec:euler}) starts at $(t_0,y_0)$ and follows the direction field from each of its corners for the time $h$. Its slopes are values of $f$, bounded by $M$, so it stays within the distance $M|t-t_0|\le b$ of $y_0$ for $|t-t_0|\le\alpha$, and it remains in $R$; this is the reason for the choice of $\alpha$. The polygons for all steps $h$ are uniformly bounded and \emph{equicontinuous}: they all satisfy $|y_h(t)-y_h(s)|\le M|t-s|$. The theorem of Arzel\`a and Ascoli then provides a sequence of steps $h_k\to0$ for which the polygons converge uniformly on $[t_0-\alpha,t_0+\alpha]$, and the uniform continuity of $f$ on $R$ shows that the limit satisfies the integral equation \eqref{eq:integral-equation}, hence the initial value problem (\cref{prop:integral-equation}). The details are in Hartman~\cite[Chap.~II]{hartman2002} and Coddington and Levinson~\cite[Chap.~1]{coddington1955}, and Volume~II proves the theorem.
\end{proof}

The interval $[t_0-\alpha,t_0+\alpha]$ is short when $M$ is large compared with $b$: a solution may leave the rectangle through its top or bottom before the time $t_0\pm a$, and the bound on the slope guarantees only that this does not happen before the time $b/M$. Peano's theorem says nothing about uniqueness, and the example $y'=\sqrt{|y|}$ shows that there may be infinitely many solutions; then different sequences of Euler polygons may converge to different solutions.

\subsection{Lipschitz Conditions}

Uniqueness requires control over how fast $f$ changes with $y$.

\begin{definition}[Lipschitz condition]
\label{def:lipschitz-y}
A function $f\colon D\to\mathbb R$, $D\subseteq\mathbb R^2$, satisfies a \term{Lipschitz condition} with respect to $y$ with the constant $L\ge0$ if
\[
|f(t,y)-f(t,z)|\le L\,|y-z|\qquad\text{for all }(t,y),(t,z)\in D .
\]
It is \term{locally Lipschitz} with respect to $y$ if every point of $D$ has a neighbourhood in $D$ on which $f$ satisfies such a condition, with a constant that may depend on the neighbourhood.
\end{definition}

This is the Lipschitz continuity of \cref{def:lipschitz} for each of the functions $y\mapsto f(t,y)$, with a constant that does not depend on $t$. The easiest way to verify it is through the derivative, as in \cref{cor:lipschitz-derivative}.

\begin{proposition}[A test for the Lipschitz condition]
\label{prop:lipschitz-test}
Let $D\subseteq\mathbb R^2$ be a set whose vertical sections $\{y:(t,y)\in D\}$ are intervals, and let $f$ be continuous on $D$ with a partial derivative $\partial f/\partial y$ that satisfies $|\partial f/\partial y|\le L$ on $D$. Then $f$ satisfies a Lipschitz condition with respect to $y$ with the constant $L$ on $D$. If $D$ is open and $\partial f/\partial y$ is continuous on $D$, then $f$ is locally Lipschitz with respect to $y$.
\end{proposition}

\begin{proof}
For fixed $t$, apply \cref{cor:lipschitz-derivative} to $y\mapsto f(t,y)$ on the interval $\{y:(t,y)\in D\}$. For the second statement, every point of the open set $D$ is the centre of a closed rectangle contained in $D$; on it the continuous function $\partial f/\partial y$ is bounded, by the extreme value theorem (\cref{thm:evt-rn}), and the first statement applies.
\end{proof}

The functions $y\sin t$ and $\sin y$ satisfy a Lipschitz condition with the constant $1$ on the whole plane. The function $t^2+y^2$ has $\partial f/\partial y=2y$, which is bounded on every bounded set but not on the plane; it is locally Lipschitz, but satisfies no Lipschitz condition on the whole plane, and the blow-up of solutions of $y'=y^2$ is connected with this. The function $\sqrt{|y|}$ is not Lipschitz near $y=0$: $|\sqrt{|y|}-0|/|y-0|=1/\sqrt{|y|}$ is unbounded as $y\to0$. Neither is $y^{2/3}$, whose derivative $2y^{-1/3}/3$ is unbounded near $0$; the equation $y'=y^{2/3}$ has, besides $y\equiv0$, the solutions $(t-c)^3/27$, and through every point of the axis it has infinitely many solutions, which are $0$ on an interval $[d,c]$ and equal to $(t-c)^3/27$ for $t>c$ and to $(t-d)^3/27$ for $t<d$.

\subsection{Gr\"onwall's Inequality and Uniqueness}

Uniqueness rests on an inequality that turns an integral inequality for a function into an explicit bound on it.

\begin{lemma}[Gr\"onwall's inequality]
\label{lem:gronwall}
Let $u$ be continuous on $[t_0,T]$, and let $c\in\mathbb R$ and $L\ge0$ be constants with
\[
u(t)\le c+L\int_{t_0}^tu(s)\,ds\qquad\text{for all }t\in[t_0,T].
\]
Then $u(t)\le c\,e^{L(t-t_0)}$ for all $t\in[t_0,T]$.
\end{lemma}

\begin{proof}
Let $U(t)\coloneqq c+L\int_{t_0}^tu(s)\,ds$. By \cref{thm:ftc1}, $U$ is differentiable with $U'=Lu\le LU$, since $u\le U$ and $L\ge0$. Hence
\[
\bigl(e^{-L(t-t_0)}U(t)\bigr)'=e^{-L(t-t_0)}\bigl(U'(t)-LU(t)\bigr)\le0 ,
\]
so $e^{-L(t-t_0)}U(t)$ does not increase (\cref{cor:monotone-derivative}) and is at most its value $U(t_0)=c$ at $t_0$. Therefore $u(t)\le U(t)\le ce^{L(t-t_0)}$.
\end{proof}

The bound is attained: $u(t)=ce^{L(t-t_0)}$ satisfies the hypothesis with equality. The differential form of the lemma, that $u'\le Lu$ implies $u(t)\le u(t_0)e^{L(t-t_0)}$, has the same proof and was used in \cref{thm:linearised-stability}.

\begin{theorem}[Uniqueness and continuous dependence]
\label{thm:uniqueness}
Let $f$ be continuous on $D\subseteq\mathbb R^2$ and satisfy a Lipschitz condition with respect to $y$ with the constant $L$ on $D$. If $y$ and $z$ are solutions of $y'=f(t,y)$ on an interval $I$ containing $t_0$, then
\begin{equation}
|y(t)-z(t)|\le|y(t_0)-z(t_0)|\,e^{L|t-t_0|}\qquad\text{for all }t\in I .
\label{eq:continuous-dependence}
\end{equation}
In particular, if $y(t_0)=z(t_0)$, then $y=z$ on $I$.
\end{theorem}

\begin{proof}
By \cref{prop:integral-equation}, for $t\in I$,
\[
y(t)-z(t)=y(t_0)-z(t_0)+\int_{t_0}^t\bigl(f(s,y(s))-f(s,z(s))\bigr)\,ds .
\]
For $t\ge t_0$, the function $u\coloneqq|y-z|$ therefore satisfies $u(t)\le u(t_0)+L\int_{t_0}^tu(s)\,ds$, by the Lipschitz condition and the triangle inequality for integrals, and \cref{lem:gronwall}, applied on $[t_0,T]$ for every $T\in I$ with $T>t_0$, gives \eqref{eq:continuous-dependence}. For $t\le t_0$, the functions $\tilde y(t)\coloneqq y(2t_0-t)$ and $\tilde z(t)\coloneqq z(2t_0-t)$ solve the equation $\tilde y'=-f(2t_0-t,\tilde y)$, whose right side satisfies a Lipschitz condition with the same constant, and the first case applies to them.
\end{proof}

The bound \eqref{eq:continuous-dependence} says more than uniqueness. Solutions that start close together stay close together on bounded intervals: an error $\varepsilon$ in the initial value causes an error of at most $\varepsilon e^{L|t-t_0|}$ at the time $t$. The bound is sharp, as the equation $y'=Ly$ shows, whose solutions differ by $(y_0-z_0)e^{L(t-t_0)}$. On long intervals the factor $e^{L|t-t_0|}$ can be enormous, and then small errors in the initial data, or in a computation, grow until the solution is effectively unpredictable; for equations with a stable equilibrium, on the other hand, differences may even decrease, as \cref{prop:linear-dependence,thm:linearised-stability} showed. Since many right sides are Lipschitz only locally, the following form is the one used most often.

\begin{corollary}[Local uniqueness]
\label{cor:local-uniqueness}
Let $D\subseteq\mathbb R^2$ be open, and let $f$ be continuous on $D$ and locally Lipschitz with respect to $y$, for instance with a continuous partial derivative $\partial f/\partial y$. If two solutions of $y'=f(t,y)$ on an interval $I$ agree at one point of $I$, they agree on all of $I$. Consequently, two integral curves with a point in common coincide over the intersection of their intervals.
\end{corollary}

\begin{proof}
Let $y$ and $z$ be the solutions, with $y(t_1)=z(t_1)$, and suppose that $y(t_2)\ne z(t_2)$ for some $t_2>t_1$ in $I$; the case $t_2<t_1$ is symmetric. Let $s\coloneqq\sup\{t\in[t_1,t_2]:y=z\text{ on }[t_1,t]\}$. By continuity $y(s)=z(s)$, so $s<t_2$. Choose a neighbourhood $Q\subseteq D$ of the point $(s,y(s))$ on which $f$ satisfies a Lipschitz condition, a small rectangle, say. By continuity both graphs lie in $Q$ for $t\in[s,s+\varepsilon]$ with some $\varepsilon>0$, and \cref{thm:uniqueness}, applied on $Q$, gives $y=z$ on $[s,s+\varepsilon]$, contradicting the definition of $s$. The last statement is the first one applied on the intersection of the two intervals, which is an interval containing the common point. A continuous $\partial f/\partial y$ makes $f$ locally Lipschitz by \cref{prop:lipschitz-test}.
\end{proof}

For the equations $y'=\sqrt{|y|}$ and $y'=y^{2/3}$, uniqueness fails only at the points of the axis $y=0$, where the Lipschitz condition fails; away from the axis, solutions are unique, and the non-uniqueness consists in the freedom of when to leave the axis.

\begin{note}{Funnels}
When solutions are not unique, they fill a region. For $y'=\sqrt{|y|}$, the solutions through $(0,0)$ have at the time $t>0$ exactly the values in $[0,t^2/4]$, bounded by the lowest solution $y\equiv0$ and the highest solution $t^2/4$, which leaves the axis immediately. A theorem of Kneser states in general that for a continuous $f$ on the rectangle of Peano's theorem and every $t_1$ in its interval $[t_0-\alpha,t_0+\alpha]$, the values $y(t_1)$ of all solutions on the interval between $t_0$ and $t_1$ form a closed bounded interval, so that the integral curves through one point sweep out a region, the \emph{funnel} of the point~\cite[Chap.~II]{hartman2002}. Uniqueness is also guaranteed by conditions weaker than the Lipschitz condition, such as Osgood's condition $|f(t,y)-f(t,z)|\le\omega(|y-z|)$ with a continuous increasing $\omega$ that is positive for $\rho>0$, with $\omega(0)=0$ and $\int_0d\rho/\omega(\rho)=\infty$; for $\omega(\rho)=\sqrt\rho$ the integral converges, in accordance with the non-uniqueness for $\sqrt{|y|}$~\cite[Chap.~III]{hartman2002}.
\end{note}

\subsection{Successive Approximations}

When $f$ is Lipschitz, the solution not only is unique, but can be constructed by iterating the integral equation.

\begin{theorem}[Picard--Lindel\"of]
\label{thm:picard-lindelof}
Let $f$ be continuous on the rectangle $R=[t_0-a,t_0+a]\times[y_0-b,y_0+b]$ and satisfy a Lipschitz condition with respect to $y$ with a constant $L>0$ on $R$, and let $M$ and $\alpha$ be as in \cref{thm:peano}. Then the initial value problem $y'=f(t,y)$, $y(t_0)=y_0$, has exactly one solution $y$ on $[t_0-\alpha,t_0+\alpha]$. The \term{Picard iterates}
\[
y_0(t)\coloneqq y_0,\qquad y_{k+1}(t)\coloneqq y_0+\int_{t_0}^tf\bigl(s,y_k(s)\bigr)\,ds\qquad(k\in\mathbb N)
\]
are defined on this interval, their graphs lie in $R$, and
\begin{gather*}
|y_{k+1}(t)-y_k(t)|\le\frac{ML^k|t-t_0|^{k+1}}{(k+1)!},\\
|y(t)-y_k(t)|\le\frac ML\sum_{j>k}\frac{(L|t-t_0|)^j}{j!} ,
\end{gather*}
so that $y_k\to y$ uniformly on $[t_0-\alpha,t_0+\alpha]$.
\end{theorem}

\begin{proof}[sketch]
Volume~II proves the theorem in full, as an application of the contraction mapping theorem. Let $|t-t_0|\le\alpha$. First, by induction, each $y_k$ is continuous with graph in $R$: if this holds for $y_k$, the integrand $f(s,y_k(s))$ is continuous and bounded by $M$, so $y_{k+1}$ is continuous and $|y_{k+1}(t)-y_0|\le M|t-t_0|\le M\alpha\le b$. Second, $|y_1(t)-y_0(t)|\le M|t-t_0|$, and if the stated bound holds for $k-1$, then for $t\ge t_0$
\begin{multline*}
|y_{k+1}(t)-y_k(t)|\le\int_{t_0}^tL\,|y_k(s)-y_{k-1}(s)|\,ds\\\le\int_{t_0}^tL\,\frac{ML^{k-1}(s-t_0)^k}{k!}\,ds=\frac{ML^k(t-t_0)^{k+1}}{(k+1)!},
\end{multline*}
and similarly for $t<t_0$. Third, the bounds are the terms of the convergent series $(M/L)(e^{L\alpha}-1)$ (\cref{thm:exp-series}), so the series $y_0+\sum_k(y_{k+1}-y_k)$ converges at every $t$, uniformly in $t$; its sum $y$ is continuous, $f(s,y_k(s))\to f(s,y(s))$ uniformly by the Lipschitz condition, and passing to the limit in the definition of $y_{k+1}$ shows that $y$ satisfies the integral equation \eqref{eq:integral-equation}, hence the initial value problem. The bound for $y-y_k$ is the tail of the series. Finally, uniqueness is \cref{thm:uniqueness}. Complete proofs are in Hartman~\cite[Chap.~II]{hartman2002}, Coddington and Levinson~\cite[Chap.~1]{coddington1955} and Teschl~\cite[Chap.~2]{teschl2012}.
\end{proof}

For $y'=y$ with $y(0)=1$ the iteration can be carried out exactly. If $y_k(t)=\sum_{j=0}^kt^j/j!$, then
\[
y_{k+1}(t)=1+\int_0^t\sum_{j=0}^k\frac{s^j}{j!}\,ds=1+\sum_{j=0}^k\frac{t^{j+1}}{(j+1)!}=\sum_{j=0}^{k+1}\frac{t^j}{j!} ,
\]
so the Picard iterates are the Taylor polynomials of $e^t$, and their convergence is that of the exponential series (\cref{thm:exp-series}). For $y'=y^2$ with $y(0)=1$, the first iterates are $1+t$ and $1+t+t^2+t^3/3$, which agree with the geometric series of the solution $1/(1-t)$ up to the terms of degree $1$ and $2$. Here $M=(1+b)^2$ on the rectangle, and $\alpha\le b/(1+b)^2\le1/4$, the maximum being attained at $b=1$; so the theorem guarantees the solution only on $[-1/4,1/4]$, although it exists on $(-\infty,1)$. The theorem is local, and the next result extends its solutions as far as possible.

\subsection{Maximal Solutions}

\begin{theorem}[Maximal solutions]
\label{thm:maximal-solution}
Let $D\subseteq\mathbb R^2$ be open, and let $f$ and $\partial f/\partial y$ be continuous on $D$. For every $(t_0,y_0)\in D$ the initial value problem $y'=f(t,y)$, $y(t_0)=y_0$, has a solution $y$ on an open interval $(\omega_-,\omega_+)$ such that every solution of the problem is a restriction of $y$; it is the \term{maximal solution}. As $t\to\omega_+$, the point $(t,y(t))$ eventually leaves every compact subset of $D$: for every compact $K\subseteq D$ there is $t_K<\omega_+$ with $(t,y(t))\notin K$ for $t_K<t<\omega_+$, and similarly as $t\to\omega_-$. In particular, if $D=J\times\mathbb R$ for an open interval $J$ with the right end $\beta$, then either $\omega_+=\beta$ or $|y(t)|\to\infty$ as $t\to\omega_+$.
\end{theorem}

\begin{proof}[sketch]
By \cref{cor:local-uniqueness}, any two solutions of the problem agree on the intersection of their intervals, so the union of their graphs is the graph of a solution on the union of their intervals, which is an interval; by \cref{thm:picard-lindelof} on small rectangles, it is open. Suppose that $(t,y(t))$ returns to a compact $K\subseteq D$ for times arbitrarily close to $\omega_+<\infty$. On a slightly larger compact set, $f$ is bounded and Lipschitz with respect to $y$, so \cref{thm:picard-lindelof} gives solutions through every point of $K$ on intervals of a common length $\alpha>0$; starting such a solution at a point $(t_1,y(t_1))\in K$ with $\omega_+-t_1<\alpha$ extends $y$ beyond $\omega_+$, a contradiction. For $D=J\times\mathbb R$ and $\omega_+<\beta$, apply this to $K=[t_0,\omega_+]\times[-r,r]$ for every $r$. The details are in Hartman~\cite[Chap.~II]{hartman2002} and Teschl~\cite[Chap.~2]{teschl2012}.
\end{proof}

The last statement is the \term{blow-up alternative}: a maximal solution on a strip ends only at the edge of the strip or by tending to infinity. A solution cannot simply stop, nor can it oscillate ever faster in a bounded region; it can do so only if the right side is unbounded there, which the hypotheses of the theorem exclude on compact sets. Together with Gr\"onwall's inequality, the alternative gives a simple criterion for solutions to exist on the whole interval.

\begin{corollary}[Global existence]
\label{cor:global-existence}
Let $J$ be an open interval, let $f$ and $\partial f/\partial y$ be continuous on $J\times\mathbb R$, and suppose that $|f(t,y)|\le A+B|y|$ for all $(t,y)$, with constants $A,B\ge0$. Then every maximal solution is defined on all of $J$.
\end{corollary}

\begin{proof}
Let $y$ be a maximal solution with $y(t_0)=y_0$, and suppose that $\omega_+$ is less than the right end of $J$. For $t_0\le t<\omega_+$, the integral equation gives
\[
|y(t)|\le|y_0|+\int_{t_0}^t\bigl(A+B|y(s)|\bigr)\,ds\le|y_0|+A(\omega_+-t_0)+B\int_{t_0}^t|y(s)|\,ds ,
\]
so by \cref{lem:gronwall} $|y(t)|\le(|y_0|+A(\omega_+-t_0))e^{B(\omega_+-t_0)}$ on every interval $[t_0,T]$ with $T<\omega_+$. Thus $y$ is bounded near $\omega_+$, contradicting \cref{thm:maximal-solution}. The same argument applies at $\omega_-$.
\end{proof}

Applied on every bounded open interval whose closure lies in the interval of continuity of the coefficients, where continuous coefficients are bounded, the corollary shows again that linear equations have solutions on the whole interval of continuity of their coefficients, as \cref{thm:linear-first-order} found by a formula. The equation $y'=\sin y$ has solutions on all of $\mathbb R$, since $|\sin y|\le1$. For $y'=y^2$ and $y'=t^2+y^2$ the right side grows quadratically, the corollary does not apply, and solutions blow up.

The results of this section up to this point hold for systems $\mathbf y'=\mathbf f(t,\mathbf y)$ with a vector-valued unknown as well~\cite[Chap.~II]{hartman2002}; Chapter~25 uses them in this form. The statements are the same once absolute values of vectors are replaced by the Euclidean norm, the interval $[y_0-b,y_0+b]$ in the rectangles by the closed ball of radius $b$ about $\mathbf y_0$, and the partial derivative $\partial f/\partial y$ by the Jacobian matrix $D_{\mathbf y}\mathbf f$ with respect to $\mathbf y$, bounded in the norm $\|\cdot\|_2$. Most proofs carry over word for word, with the triangle inequality for integrals of vector-valued functions (\cref{prop:mvt-integral-rn}) in place of the scalar one. The exception is the test of \cref{prop:lipschitz-test}: a vector-valued function satisfies no mean value theorem, and the proof uses the mean value inequality (\cref{thm:mean-value-inequality}) instead, for sets whose sections $\{\mathbf y:(t,\mathbf y)\in D\}$ are convex. The comparison theorem below has no such simple extension, because vectors are not ordered.

\subsection{Comparison}

Solutions of an equation that cannot be solved can often be estimated by solutions of one that can. The tool is the following comparison principle.

\begin{theorem}[Comparison]
\label{thm:comparison}
Let $J$ be an open interval, let $f$ and $\partial f/\partial y$ be continuous on $J\times\mathbb R$, let $[t_0,T]\subseteq J$, and let $y$ be a solution of $y'=f(t,y)$ on $[t_0,T]$. If $z$ is differentiable on $[t_0,T]$ with
\[
z'(t)\le f\bigl(t,z(t)\bigr)\quad\text{for }t\in[t_0,T],\qquad z(t_0)\le y(t_0),
\]
then $z(t)\le y(t)$ for all $t\in[t_0,T]$. If instead $z'\ge f(t,z)$ and $z(t_0)\ge y(t_0)$, then $z\ge y$ on $[t_0,T]$.
\end{theorem}

\begin{proof}
Let $w\coloneqq z-y$ and suppose that $w(t_2)>0$ for some $t_2\in(t_0,T]$. Let $t_1\coloneqq\sup\{t\in[t_0,t_2]:w(t)\le0\}$. By continuity $w(t_1)\le0$, and $w>0$ on $(t_1,t_2]$, so $w(t_1)=0$ and $t_1<t_2$. The set $Q$ of points $(t,\eta)$ with $t_1\le t\le t_2$ and $\eta$ between $y(t)$ and $z(t)$ is closed and bounded, so $|\partial f/\partial y|$ has a maximum $L$ on it (\cref{thm:evt-rn}), and by the mean value theorem $f(t,z)-f(t,y)\le L(z-y)$ for $t\in[t_1,t_2]$. Hence on $[t_1,t_2]$
\[
w'=z'-y'\le f(t,z)-f(t,y)\le Lw ,
\]
so $(e^{-Lt}w)'\le0$, and $e^{-Lt_2}w(t_2)\le e^{-Lt_1}w(t_1)=0$, a contradiction. The second statement follows by applying the first to $-z$ and $-y$, which satisfy the comparison for the right side $-f(t,-y)$.
\end{proof}

A function with $z'\le f(t,z)$ is called a \term{subsolution}, and one with $z'\ge f(t,z)$ a \term{supersolution}. The solution of $y'=t^2+y^2$ with $y(0)=0$, which cannot be written down, shows how the theorem is used; let $(\omega_-,\omega_+)$ be its maximal interval. Since $y'\ge0$, $y$ increases and $y\ge0$ for $t\ge0$. On $[0,1]$, $f(t,z)=t^2+z^2\le1+z^2$, so $z(t)=\tan t$, which satisfies $z'=1+z^2$, is a supersolution with $z(0)=0$, and $0\le y(t)\le\tan t\le\tan1$ for $0\le t<\min(1,\omega_+)$. By the blow-up alternative, $y$ cannot end while it is bounded, so $\omega_+>1$. Next, $y'\ge t^2$ gives $y(1)\ge1/3$, and for $t\ge1$, $f(t,z)\ge1+z^2$, so $z(t)=\tan(t-1+\arctan(1/3))$ is a subsolution with $z(1)=1/3\le y(1)$. It tends to $\infty$ as $t$ approaches $t^*\coloneqq1+\pi/2-\arctan(1/3)=2.24905\ldots$, and $y\ge z$ as long as both exist; since a continuous $y$ cannot tend to $\infty$ at an interior point of its interval, $\omega_+\le t^*$. So the solution blows up at a time $\omega_+$ with $1<\omega_+\le t^*<2.2491$. A numerical computation with the classical fourth-order Runge--Kutta method~\cite[Chap.~II]{hairer1993}, applied to $y$ up to the time at which $y$ reaches $1$ and then to the reciprocal $w=1/y$, which satisfies $w'=-(t^2w^2+1)$ and crosses $0$ at the blow-up time, gave $\omega_+=2.003147359$ for each of the steps $h=0.01$, $0.005$ and $0.0025$; the three results agree in all the digits shown. This also matches the steep rise of the curve through the origin in \cref{fig:direction-fields}C. Since $-y(-t)$ is a solution with the same initial value, it equals $y$ by uniqueness, so $y$ is odd and $\omega_-=-\omega_+$.

\section{Euler's Method: A First Look at Numerical Solution}
\label{sec:euler}

\subsection{Following the Direction Field}

When no formula is available, a solution can be computed approximately at a sequence of times $t_k=t_0+kh$ with a \term{step} $h>0$. Over one step, the integral equation \eqref{eq:integral-equation} gives
\[
y(t_{k+1})=y(t_k)+\int_{t_k}^{t_{k+1}}f\bigl(s,y(s)\bigr)\,ds .
\]
Replacing the integrand by its value at the left end of the interval, as in the left rectangle rule of \cref{def:quadrature-rules}, and the unknown $y(t_k)$ by an approximation $y_k$, gives \term{Euler's method}
\begin{equation}
y_{k+1}=y_k+h\,f(t_k,y_k)\qquad(k=0,1,2,\ldots).
\label{eq:euler-method}
\end{equation}
Geometrically, the method walks along the direction field: from the point $(t_k,y_k)$ it follows the segment of slope $f(t_k,y_k)$ for the time $h$, and there it takes a new direction (\cref{fig:euler}A). The broken line through the points $(t_k,y_k)$ is the \term{Euler polygon}, the curve of the proof of Peano's theorem. \Cref{alg:euler} states the method. The computations of this section were done in plain Python on the machine described in \cref{sec:systems-practice}, in binary64 arithmetic unless stated otherwise; exact reference values come from the formulas for the solutions.

\begin{algorithm}
\begin{algorithmic}[1]
\Require the right side $f$; the initial point $(t_0,y_0)$; the step $h>0$; the number of steps $N$
\Ensure approximations $y_k$ of $y(t_k)$ at $t_k=t_0+kh$, $k=0,\ldots,N$
\Name{Euler}{f,t_0,y_0,h,N}
\For{$k=0,1,\ldots,N-1$}
  \State $y_{k+1}\gets y_k+h\,f(t_k,y_k)$ \Comment{follow the direction field for the time $h$}
  \State $t_{k+1}\gets t_0+(k+1)h$
\EndFor
\State \Return $y_0,\ldots,y_N$
\end{algorithmic}
\caption{Euler's method for the initial value problem $y'=f(t,y)$, $y(t_0)=y_0$. Each step costs one evaluation of $f$. The times are computed as $t_0+(k+1)h$, which avoids the accumulation of rounding errors that repeated addition of $h$ would cause.\label{alg:euler}}
\end{algorithm}

For $y'=y$ with $y(0)=1$, the method gives $y_{k+1}=(1+h)y_k$, so $y_k=(1+h)^k$. With $h=1/n$, the approximation of $y(1)=e$ after $n$ steps is $(1+1/n)^n$, the compound interest of \cref{def:e}, and \cref{thm:exp-limit} says that it converges to $e$; \cref{thm:exp-limit-matrix} made the same observation for systems $\mathbf x'=A\mathbf x$. Euler's method thus generalises the limit that defined $e$, and its convergence for every equation with a Lipschitz right side is the subject of the next subsection.

\subsection{Local and Global Errors}

The exact solution does not satisfy the recurrence \eqref{eq:euler-method} exactly. The amount by which it fails,
\[
\tau_{k+1}\coloneqq y(t_{k+1})-y(t_k)-h\,f\bigl(t_k,y(t_k)\bigr)=y(t_{k+1})-y(t_k)-h\,y'(t_k),
\]
is the \term{local truncation error} of the step. If $y$ is twice continuously differentiable, Taylor's theorem with the Lagrange remainder (\cref{thm:taylor-lagrange}) gives $\tau_{k+1}=h^2y''(\xi_k)/2$ with some $\xi_k\in(t_k,t_{k+1})$, so
\[
|\tau_{k+1}|\le\frac{M_2h^2}2,\qquad M_2\coloneqq\max|y''| .
\]
The solution is twice continuously differentiable when $f$ is continuously differentiable, since then $y''=f_t(t,y)+f_y(t,y)\,f(t,y)$ by the chain rule (\cref{thm:chain-rule-rn}). The \term{global error} $e_k\coloneqq y_k-y(t_k)$ accumulates the local errors of all earlier steps, each of them propagated through the later steps. The following discrete form of Gr\"onwall's inequality controls the propagation.

\begin{lemma}[Discrete Gr\"onwall inequality]
\label{lem:discrete-gronwall}
Let $h>0$, $L>0$ and $\delta\ge0$, and let the numbers $\varepsilon_0,\varepsilon_1,\ldots\ge0$ satisfy $\varepsilon_{k+1}\le(1+hL)\varepsilon_k+\delta$ for all $k$. Then
\[
\varepsilon_k\le e^{khL}\varepsilon_0+\frac\delta{hL}\bigl(e^{khL}-1\bigr)\qquad(k\in\mathbb N).
\]
\end{lemma}

\begin{proof}
By induction on $k$, $\varepsilon_k\le(1+hL)^k\varepsilon_0+\delta\sum_{j=0}^{k-1}(1+hL)^j=(1+hL)^k\varepsilon_0+\delta\bigl((1+hL)^k-1\bigr)/(hL)$, by the formula for a geometric sum. Since $1+hL\le e^{hL}$ by \cref{prop:exp-tangent}, the claim follows.
\end{proof}

\begin{theorem}[Convergence of Euler's method]
\label{thm:euler-convergence}
Let $f$ be continuous on $[t_0,T]\times\mathbb R$ and satisfy a Lipschitz condition with respect to $y$ with a constant $L>0$ there, let $y$ be the solution of $y'=f(t,y)$, $y(t_0)=y_0$, on $[t_0,T]$, and suppose that $y$ is twice continuously differentiable with $|y''|\le M_2$. Let $h=(T-t_0)/N$, and let $y_0^h,y_1^h,\ldots,y_N^h$ be computed by \eqref{eq:euler-method} from an initial value $y_0^h$. Then
\[
|y_k^h-y(t_k)|\le e^{L(t_k-t_0)}|y_0^h-y_0|+\frac{M_2h}{2L}\bigl(e^{L(t_k-t_0)}-1\bigr)\qquad(k=0,\ldots,N).
\]
In particular, if $y_0^h=y_0$, the error is at most $Ch$ with $C\coloneqq M_2(e^{L(T-t_0)}-1)/(2L)$, and Euler's method converges with order $1$.
\end{theorem}

\begin{proof}
Subtract the definition of $\tau_{k+1}$ from \eqref{eq:euler-method}: the error $e_k\coloneqq y_k^h-y(t_k)$ satisfies
\[
e_{k+1}=e_k+h\bigl(f(t_k,y_k^h)-f(t_k,y(t_k))\bigr)-\tau_{k+1} .
\]
By the Lipschitz condition and the bound on $\tau_{k+1}$, $|e_{k+1}|\le(1+hL)|e_k|+M_2h^2/2$, and \cref{lem:discrete-gronwall} with $\delta=M_2h^2/2$ and $kh=t_k-t_0$ gives the claim.
\end{proof}

The proof has a structure that recurs in the analysis of every numerical method for differential equations: a local error of order $h^2$ per step, made in each of $(T-t_0)/h$ steps, and a stability estimate, here the Lipschitz condition through Gr\"onwall's inequality, which ensures that the errors are amplified by at most the bounded factor $e^{L(T-t_0)}$. The result is a global error of order $h$. The bound is usually pessimistic, since $e^{L(T-t_0)}$ describes the worst possible growth of differences, as in \cref{thm:uniqueness}. The hypothesis of a Lipschitz condition on the whole strip can be weakened to a neighbourhood of the solution, at the price of requiring $h$ to be small enough that the Euler polygon stays in it~\cite[Chap.~II]{hairer1993}.

\subsection{Experiments}

For $y'=y$ on $[0,1]$, the theorem with $L=1$ and $M_2=e$ bounds the error of $(1+1/n)^n$ by $e(e-1)h/2=2.34h$, with $h=1/n$. The actual error is close to $eh/2=1.36h$. Indeed, by the logarithmic series (\cref{thm:log-series}), $\ln(1+h)/h=1-h/2+O(h^2)$, so
\[
(1+h)^{1/h}=\exp\Bigl(1-\frac h2+O(h^2)\Bigr)=e\Bigl(1-\frac h2+O(h^2)\Bigr),\qquad e-(1+h)^{1/h}=\frac{eh}2+O(h^2).
\]
\Cref{tab:euler-exp} shows the errors: each halving of $h$ halves the error, and $n$ times the error approaches $e/2=1.3591$.

Better methods sample the direction field more than once per step. The \term{improved Euler method}, or \term{method of Heun}, takes a trial Euler step and averages the slopes at its two ends, and the \term{midpoint method} uses the slope at the middle of a half step:
\begin{align*}
y_{k+1}&=y_k+\frac h2\Bigl(f(t_k,y_k)+f\bigl(t_{k+1},y_k+hf(t_k,y_k)\bigr)\Bigr)&&\text{(Heun)},\\
y_{k+1}&=y_k+h\,f\Bigl(t_k+\frac h2,\,y_k+\frac h2f(t_k,y_k)\Bigr)&&\text{(midpoint)}.
\end{align*}
Each step costs two evaluations of $f$. They correspond to the trapezoidal rule and the midpoint rule of \cref{sec:quadrature}, with the unknown value at the right end or in the middle replaced by an Euler prediction. For $y'=y$ both reduce to $y_{k+1}=(1+h+h^2/2)y_k$, the Taylor polynomial of degree $2$ of $e^h$, and the same expansion as above, now with $\ln(1+h+h^2/2)=h-h^3/6+O(h^4)$, gives the error $eh^2/6+O(h^3)$. The last columns of \cref{tab:euler-exp} confirm it: each halving of $h$ divides the error by $4$, and $n^2$ times the error approaches $e/6=0.4530$. A method whose global error behaves like $Ch^p$ is said to have \term{order} $p$; Heun's method and the midpoint method have order $2$.

\begin{table}
\sdfloat{%
\begin{sdtable}{colspec={rllrllr}, row{1}={font=\small\bfseries}}
\toprule
$n$ & Euler error & ratio & $n\cdot$error & Heun error & ratio & $n^2\cdot$error \\
\midrule
$10$ & $1.245\times10^{-1}$ &  & $1.2454$ & $4.201\times10^{-3}$ &  & $0.4201$ \\
$20$ & $6.498\times10^{-2}$ & 1.92 & $1.2997$ & $1.091\times10^{-3}$ & 3.85 & $0.4363$ \\
$40$ & $3.322\times10^{-2}$ & 1.96 & $1.3287$ & $2.779\times10^{-4}$ & 3.93 & $0.4446$ \\
$80$ & $1.680\times10^{-2}$ & 1.98 & $1.3438$ & $7.013\times10^{-5}$ & 3.96 & $0.4488$ \\
$160$ & $8.446\times10^{-3}$ & 1.99 & $1.3514$ & $1.761\times10^{-5}$ & 3.98 & $0.4509$ \\
$320$ & $4.235\times10^{-3}$ & 1.99 & $1.3553$ & $4.414\times10^{-6}$ & 3.99 & $0.4520$ \\
$640$ & $2.121\times10^{-3}$ & 2.00 & $1.3572$ & $1.105\times10^{-6}$ & 4.00 & $0.4525$ \\
$1280$ & $1.061\times10^{-3}$ & 2.00 & $1.3582$ & $2.764\times10^{-7}$ & 4.00 & $0.4528$ \\
\bottomrule
\end{sdtable}}{Errors $e-y_n$ at $t=1$ of Euler's method and of Heun's method for $y'=y$, $y(0)=1$, with $n$ steps of length $h=1/n$; the midpoint method gives the same values as Heun's method for this equation. The ratios are those of consecutive errors.\label{tab:euler-exp}}
\end{table}

For a nonlinear equation the two second-order methods differ. \Cref{tab:euler-logistic} shows the errors at $t=4$ for the logistic equation $y'=y(1-y)$ with $y(0)=0.1$, whose solution is $1/(1+9e^{-t})$ by \eqref{eq:logistic-solution}. The ratios of consecutive errors approach $2$ for Euler's method and $4$ for the two others, and for this equation the midpoint method is about three times as accurate as Heun's method. The comparison at equal cost favours the methods of order $2$: with $n=80$ steps, Heun's method uses $160$ evaluations of $f$ and has the error $8.7\times10^{-5}$, while Euler's method with $160$ evaluations has the error $4.6\times10^{-4}$, and the gap widens as the accuracy demanded grows. Methods of higher order, the Runge--Kutta methods and the multistep methods, and the techniques that adjust the step to the solution, are described by Hairer, N{\o}rsett and Wanner~\cite{hairer1993} and by Iserles~\cite{iserles2008}.

\begin{table}
\sdfloat{%
\begin{sdtable}{colspec={rrllllll}, colsep=4pt, row{1}={font=\small\bfseries}}
\toprule
$n$ & $h$ & Euler & ratio & Heun & ratio & midpoint & ratio \\
\midrule
$10$ & $0.4$ & $9.909\times10^{-3}$ &  & $5.500\times10^{-3}$ &  & $1.771\times10^{-3}$ &  \\
$20$ & $0.2$ & $4.225\times10^{-3}$ & 2.35 & $1.378\times10^{-3}$ & 3.99 & $4.587\times10^{-4}$ & 3.86 \\
$40$ & $0.1$ & $1.960\times10^{-3}$ & 2.16 & $3.462\times10^{-4}$ & 3.98 & $1.167\times10^{-4}$ & 3.93 \\
$80$ & $0.05$ & $9.447\times10^{-4}$ & 2.07 & $8.689\times10^{-5}$ & 3.98 & $2.944\times10^{-5}$ & 3.96 \\
$160$ & $0.025$ & $4.639\times10^{-4}$ & 2.04 & $2.177\times10^{-5}$ & 3.99 & $7.392\times10^{-6}$ & 3.98 \\
$320$ & $0.0125$ & $2.298\times10^{-4}$ & 2.02 & $5.449\times10^{-6}$ & 4.00 & $1.852\times10^{-6}$ & 3.99 \\
$640$ & $0.00625$ & $1.144\times10^{-4}$ & 2.01 & $1.363\times10^{-6}$ & 4.00 & $4.636\times10^{-7}$ & 4.00 \\
\bottomrule
\end{sdtable}}{Errors $y(4)-y_n$ of three methods for the logistic equation $y'=y(1-y)$, $y(0)=0.1$, with $n$ steps of length $h=4/n$; all errors are positive.\label{tab:euler-logistic}}
\end{table}

\subsection{Stability and Stiffness}

Convergence as $h\to0$ does not mean that every step is acceptable. Apply Euler's method to the \term{test equation} $y'=\lambda y$ with a constant $\lambda<0$, whose solutions decay. The method gives $y_k=(1+h\lambda)^ky_0$, which decays if and only if $|1+h\lambda|<1$, that is, $h<2/|\lambda|$. For larger steps the approximations grow in modulus with alternating signs while the solution decays: for $\lambda=-50$ and $h=0.05$, the factor is $1-50h=-1.5$, and after $20$ steps $y_{20}=(-1.5)^{20}=3325$, while $y(1)=e^{-50}=1.9\times10^{-22}$. The method is \term{stable} for this equation only when $h<2/|\lambda|=0.04$. The recurrence $y_{k+1}=(1+h\lambda)y_k$ is a difference equation, and the stability of such recurrences is a theme of Chapter~27.

The restriction matters even when the solution itself varies slowly. The equation
\[
y'=-50\,(y-\cos t),\qquad y(0)=0 ,
\]
is linear, and by \cref{sec:linear-first-order} its solution is
\[
y(t)=\frac{50}{2501}\,(50\cos t+\sin t)-\frac{2500}{2501}\,e^{-50t} .
\]
After a short transient of duration about $0.1$, the solution stays close to $\cos t$ and changes on the time scale $1$; a step of $0.1$ would follow it accurately. But Euler's method needs $h<0.04$ to remain stable, because the perturbations of the slowly varying solution decay like $e^{-50t}$ and are subject to the same condition as the test equation. With $h=0.05$ the approximations oscillate with growing amplitude, reaching $8.57$ at $t=0.25$ and $-3325$ at $t=1$; with $h=0.035$ they oscillate about the solution and settle, with an error below $0.014$ for $t\ge0.5$ (\cref{fig:euler}B). Equations whose solutions contain rapidly decaying components next to slowly varying ones are called \term{stiff}, a term introduced by Curtiss and Hirschfelder~\cite{curtiss1952}. For them, explicit methods such as Euler's must take steps adapted to the fastest time scale long after the fast components have died out.

The remedy is to evaluate the slope at the end of the step. The \term{backward Euler method}
\[
y_{k+1}=y_k+h\,f(t_{k+1},y_{k+1})
\]
is \term{implicit}: the new value appears on both sides, and in general each step requires the solution of an equation, for instance by Newton's method (\cref{sec:newton}). For the test equation, $y_{k+1}=y_k/(1-h\lambda)$, which decays for every $h>0$ when $\lambda<0$. For the stiff equation above, the backward Euler method with $h=0.1$ is stable; its largest error, $0.16$, occurs in the first step, where the transient is not resolved, and at $t=1$ its error is $6\times10^{-4}$ (\cref{fig:euler}B). Hairer and Wanner~\cite{hairer1996stiff} treat stiff equations and the implicit methods designed for them.

\begin{figure}
\sdpanel{A}{\begin{tikzpicture}
\begin{axis}[sdplot, width=0.47\textwidth, height=0.42\textwidth, xmin=0, xmax=1.05, ymin=0.9, ymax=2.8,
  xtick={0,0.25,0.5,0.75,1}, ytick={1,1.5,2,2.5}, xlabel={$t$}, ylabel={$y$}]
\addplot[black, line width=0.9pt, domain=0:1, samples=60] {exp(x)};
\addplot[black, dashed, line width=0.6pt, mark=*, mark size=1.2pt, mark options={solid}] coordinates {(0,1) (0.5,1.5) (1,2.25)};
\addplot[black!55, line width=0.6pt, mark=*, mark size=1.2pt] coordinates {(0,1) (0.25,1.25) (0.5,1.5625) (0.75,1.953125) (1,2.44140625)};
\addplot[black, densely dotted, line width=0.7pt, mark=*, mark size=1.0pt, mark options={solid}] coordinates {(0,1) (0.125,1.125) (0.25,1.265625) (0.375,1.423828) (0.5,1.601807) (0.625,1.802032) (0.75,2.027287) (0.875,2.280697) (1,2.565785)};
\end{axis}
\end{tikzpicture}}\hfill
\sdpanel{B}{\begin{tikzpicture}
\begin{axis}[sdplot, width=0.47\textwidth, height=0.42\textwidth, xmin=0, xmax=0.72, ymin=-1.5, ymax=2.6,
  xtick={0,0.2,0.4,0.6}, ytick={-1,0,1,2}, xlabel={$t$}, ylabel={$y$}, clip=true]
\addplot[black, line width=0.9pt, domain=0:0.72, samples=200] {(50/2501)*(50*cos(deg(x))+sin(deg(x)))-(2500/2501)*exp(-50*x)};
\addplot[black!50, line width=0.5pt] coordinates {(0.000,0.0000) (0.050,2.5000) (0.100,-1.2531) (0.150,4.3672)};
\addplot[black, densely dotted, line width=0.6pt] coordinates {(0.000,0.0000) (0.035,1.7500) (0.070,0.4364) (0.105,1.4184) (0.140,0.6766) (0.175,1.2255) (0.210,0.8042) (0.245,1.1084) (0.280,0.8664) (0.315,1.0320) (0.350,0.8899) (0.385,0.9765) (0.420,0.8895) (0.455,0.9308) (0.490,0.8739) (0.525,0.8887) (0.560,0.8478) (0.595,0.8468) (0.630,0.8141) (0.665,0.8034) (0.700,0.7745)};
\addplot[only marks, mark=o, mark size=1.6pt, black] coordinates {(0.000,0.0000) (0.100,0.8292) (0.200,0.9549) (0.300,0.9553) (0.400,0.9268) (0.500,0.8858) (0.600,0.8354) (0.700,0.7766)};
\end{axis}
\end{tikzpicture}}
\caption{Euler's method. (A) Euler polygons for $y'=y$, $y(0)=1$, with the steps $h=1/2$ (dashed), $1/4$ (grey) and $1/8$ (dotted), below the solution $e^t$ (solid); their end values are $2.25$, $2.441$ and $2.566$, against $e=2.718$. (B) The stiff equation $y'=-50(y-\cos t)$, $y(0)=0$: the solution (solid), Euler's method with $h=0.05$ (grey), which leaves the window after three steps, Euler's method with $h=0.035$ (dotted), which oscillates and settles, and the backward Euler method with $h=0.1$ (circles), shown up to $t=0.7$.\label{fig:euler}}
\end{figure}

\subsection{Rounding Errors}

\Cref{thm:euler-convergence} assumes exact arithmetic. In floating-point arithmetic each step adds a rounding error, and a smaller step means more steps and more rounding errors. So the total error cannot be made arbitrarily small by decreasing $h$, and there is a smallest useful step. How small it is depends on the precision. \Cref{tab:euler-rounding} shows Euler's method for $y'=y$ on $[0,1]$ with $n=10^k$ steps, the step being the floating-point number nearest to $1/n$, computed in binary64 and in binary32 arithmetic, the latter with the 32-bit floating-point type of the NumPy library (version 2.5.5). The value $E_n=(1+1/n)^n$ that exact arithmetic would produce was computed to $50$ digits with Python's \code{decimal} module. In binary64, the computed values differed from $E_n$ by at most $9.0\times10^{-13}$ for each of these $n$ up to $10^8$, so the error is the truncation error $e-E_n$ to the digits shown. In binary32, with the unit roundoff $2^{-24}=6.0\times10^{-8}$, the error decreased to $4.9\times10^{-5}$ at $n=10^5$ and then grew; at $n=10^8$ the computed values never changed from $1$, because $hy_k=10^{-8}$ is less than half the spacing $2^{-23}$ of binary32 numbers just above $1$, so that $y_k+hy_k$ rounds to $y_k$. Volume~V analyses how rounding errors accumulate in long recurrences and sums.

\begin{table}
\sdfloat{%
\begin{sdtable}{colspec={rlll}, row{1}={font=\small\bfseries}}
\toprule
$n$ & $e-E_n$ & binary64: $|y_n-E_n|$ & binary32: $e-y_n$ \\
\midrule
$10^1$ & $1.245\times10^{-1}$ & $9.1\times10^{-17}$ & $1.245\times10^{-1}$ \\
$10^2$ & $1.347\times10^{-2}$ & $2.8\times10^{-16}$ & $1.347\times10^{-2}$ \\
$10^3$ & $1.358\times10^{-3}$ & $3.6\times10^{-15}$ & $1.361\times10^{-3}$ \\
$10^4$ & $1.359\times10^{-4}$ & $1.7\times10^{-15}$ & $1.384\times10^{-4}$ \\
$10^5$ & $1.359\times10^{-5}$ & $5.6\times10^{-15}$ & $4.944\times10^{-5}$ \\
$10^6$ & $1.359\times10^{-6}$ & $9.5\times10^{-14}$ & $1.963\times10^{-4}$ \\
$10^7$ & $1.359\times10^{-7}$ & $9.0\times10^{-13}$ & $1.222\times10^{-1}$ \\
$10^8$ & $1.359\times10^{-8}$ & $8.5\times10^{-13}$ & $1.718$ \\
\bottomrule
\end{sdtable}}{Euler's method for $y'=y$, $y(0)=1$, at $t=1$ with $n$ steps: the truncation error $e-E_n$ of exact arithmetic, the rounding error of the binary64 computation, and the total error of the binary32 computation.\label{tab:euler-rounding}}
\end{table}

\bigskip
A differential equation of the first order prescribes a slope at every point, and its solutions are the curves that follow this direction field; one initial value singles out one solution, provided the right side is Lipschitz with respect to $y$, and otherwise several solutions may share an initial point, as for $y'=\sqrt{|y|}$ and the draining tank. Separable equations are solved by integrating $dy/h(y)=g(t)\,dt$, with the zeros of $h$ as constant solutions; this gives exponential growth and decay with their doubling times and half-lives, Newton's law of cooling, the logistic curve, the blow-up of $y'=y^2$, the velocity $v_{\mathrm T}\tanh(gt/v_{\mathrm T})$ of a falling body and the catenary $a\cosh(x/a)$. Linear equations are solved by the integrating factor $e^{\int p}$ on the whole interval of continuity of their coefficients; their solutions are a particular solution plus the one-dimensional kernel of $D+p$, they respond to periodic inputs with a transient and a steady state, and they forget their initial values when $p$ is bounded below by a positive constant. Exact equations are solved by the potential of a conservative field, whose level curves are the solutions, and integrating factors, substitutions and scaling reduce further equations to these types. Autonomous equations are described by their phase line: solutions move monotonically between equilibria, the sign of $f'$ at an equilibrium decides its stability, and a parameter can make equilibria merge and vanish. Peano's theorem gives existence for continuous right sides, Gr\"onwall's inequality gives uniqueness and continuous dependence for Lipschitz ones, successive approximations construct the solution, maximal solutions end only at the boundary or by blowing up, and comparison with solvable equations locates the blow-up. Euler's method converges with order one, needs small steps for stiff equations, and meets a limit set by rounding errors.

The next chapter turns to linear equations of higher order, $y^{(k)}+a_{k-1}(t)y^{(k-1)}+\cdots+a_0(t)y=q(t)$. Their solutions again form a particular solution plus a space of solutions of the homogeneous equation, now of dimension $k$, as \cref{thm:linear-homogeneous} foreshadowed for $k=1$; for constant coefficients they are built from the functions $t^je^{\lambda t}$ for the roots $\lambda$ of a polynomial, and the forced oscillations of the second order show the resonance that \cref{sec:linear-first-order} could not show. Chapter~25 extends the phase line to the phase plane, and the linearised stability of \cref{thm:linearised-stability} to systems, where eigenvalues take the place of $f'(y^*)$.

\chapternotes

The first differential equations were solved in the decades after the invention of the calculus by Newton and Leibniz, and the methods of separation of variables and of integrating factors, the linear equation and the substitution that solves Bernoulli's equation, all belong to that period; Kline~\cite{kline1972} tells this history. The catenary was found in 1691 by Leibniz, Huygens and Johann Bernoulli~\cite{maor1994}. Verhulst proposed the logistic law of population growth in 1838~\cite{verhulst1838}, and Liouville proved in 1841 that the Riccati equation $y'=y^2+t^m$ is solvable by quadratures only in the exceptional cases recorded in the note of \cref{sec:direction-fields}~\cite{liouville1841,watson1944}. Euler introduced his method in his textbook on integral calculus, and Cauchy used Euler polygons to prove the existence of solutions; Hairer, N{\o}rsett and Wanner describe this history~\cite[Chaps.~I and~II]{hairer1993}. Peano's proof of existence for continuous right sides, in its general form for systems, appeared in 1890~\cite{peano1890}. Picard developed the method of successive approximations in his memoir of 1890~\cite{picard1890}, and in 1894 Lindel\"of presented it under a Lipschitz condition alone, without assuming that the right side has derivatives, together with a modification of his own and a simple proof of uniqueness~\cite{lindelof1894}. Gr\"onwall's inequality goes back to his paper of 1919 on the dependence of solutions on parameters~\cite{gronwall1919}. Curtiss and Hirschfelder introduced the notion of stiffness in 1952~\cite{curtiss1952}. The half-life of carbon-14 used in \cref{sec:separable} is the value of Godwin~\cite{godwin1962}.

Boyce, DiPrima and Meade~\cite{boyce2017} treat the methods of this chapter with many applications, including the population models of \cref{sec:autonomous} and the elementary theory of existence and uniqueness. Strogatz~\cite{strogatz2015} develops the qualitative theory of autonomous equations and their bifurcations, beginning with the phase line. For the theory of existence, uniqueness, maximal solutions, differential inequalities and dependence on initial values, Coddington and Levinson~\cite{coddington1955}, Hartman~\cite{hartman2002} and Teschl~\cite{teschl2012} are standard; Hartman also treats Kneser's theorem and Osgood's condition. The numerical solution of ordinary differential equations is developed by Hairer, N{\o}rsett and Wanner~\cite{hairer1993} for nonstiff and by Hairer and Wanner~\cite{hairer1996stiff} for stiff problems, and Iserles~\cite{iserles2008} gives an introduction to both.

\chapterreferences
